\documentclass[12pt, a4paper]{article}

\usepackage{amsmath, amssymb, amsthm}
\usepackage{mathtools}
\usepackage{geometry}
\usepackage{microtype}
\usepackage{setspace}
\usepackage{hyperref}
\usepackage{cleveref}
\usepackage{booktabs}
\usepackage{array}
\usepackage{xcolor}
\usepackage{tocloft}
\usepackage{titlesec}
\usepackage{fancyhdr}
\usepackage{enumitem}
\usepackage{mdframed}
\usepackage{graphicx}
\geometry{
  top=1.1in, bottom=1.1in,
  left=1.25in, right=1.25in,
  headheight=15pt
}

\onehalfspacing

\hypersetup{
  colorlinks=true,
  linkcolor=black!60!blue,
  citecolor=black!50!green,
  urlcolor=black!50!blue,
  pdftitle={Activation Manifolds as Admissibility Fields},
  pdfauthor={Flyxion}
}

\pagestyle{fancy}
\fancyhf{}
\fancyhead[L]{\small\textit{Activation Manifolds as Admissibility Fields}}
\fancyhead[R]{\small Flyxion}
\fancyfoot[C]{\thepage}
\renewcommand{\headrulewidth}{0.3pt}

\titleformat{\section}{\large\bfseries}{\thesection.}{0.6em}{}
\titleformat{\subsection}{\normalsize\bfseries}{\thesubsection.}{0.5em}{}
\titleformat{\subsubsection}{\normalsize\itshape}{\thesubsubsection.}{0.5em}{}

\theoremstyle{plain}
\newtheorem{theorem}{Theorem}[section]
\newtheorem{proposition}[theorem]{Proposition}
\newtheorem{lemma}[theorem]{Lemma}
\newtheorem{corollary}[theorem]{Corollary}

\theoremstyle{definition}
\newtheorem{definition}[theorem]{Definition}
\newtheorem{example}[theorem]{Example}

\theoremstyle{remark}
\newtheorem{remark}[theorem]{Remark}
\newtheorem{observation}[theorem]{Observation}

\newmdenv[
  linewidth=0.6pt,
  linecolor=black!40,
  backgroundcolor=black!3,
  innertopmargin=8pt,
  innerbottommargin=8pt,
  innerleftmargin=12pt,
  innerrightmargin=12pt,
  skipabove=10pt,
  skipbelow=10pt
]{principleBox}

\newcommand{\R}{\mathbb{R}}
\newcommand{\Sph}{\mathbb{S}}
\newcommand{\cM}{\mathcal{M}}
\newcommand{\cF}{\mathcal{F}}
\newcommand{\cA}{\mathcal{A}}
\newcommand{\cL}{\mathcal{L}}
\newcommand{\cG}{\mathcal{G}}
\newcommand{\cI}{\mathcal{I}}
\newcommand{\cB}{\mathcal{B}}
\newcommand{\fA}{A}
\newcommand{\norm}[1]{\left\|#1\right\|}
\newcommand{\ip}[2]{\left\langle #1,\, #2 \right\rangle}
\newcommand{\hh}{\hat{h}}
\newcommand{\ra}{\rightarrow}
\newcommand{\Ra}{\Rightarrow}
\newcommand{\lra}{\longrightarrow}
\newcommand{\simR}{\sim_R}
\newcommand{\reachset}[1]{R(#1)}
\newcommand{\WTwo}{W_2}
\newcommand{\Adm}{\mathcal{A}_{\mathrm{dm}}}
\newcommand{\ERF}{\mathrm{ERF}}
\DeclareMathOperator{\kur}{\kappa}

\begin{document}

\begin{titlepage}
\centering
\vspace*{1.5cm}

{\LARGE\bfseries Activation Manifolds as Admissibility Fields:\\[0.4em]
Reachability Geometry in Large Language Models}

\vspace{2cm}

{\large Flyxion}\\[0.3em]
{\normalsize Independent Research}

\vspace{1cm}

{\normalsize June 2026}

\vspace{2.5cm}

\begin{abstract}
\noindent
Three independent empirical results from June 2026 --- normalization-revealed
activation geometry \cite{liu2026normgeometry}, ICA-recovered non-Gaussian
directional structure \cite{liu2026icalens}, and trajectory-field
self-distillation in world models \cite{stapf2026wmsd} --- converge on a
single geometric picture that has not previously been articulated in a unified
form. We argue that what these three programs are collectively discovering in
activation space is not feature geography but an emergent shadow of a deeper
constraint topology, which we formalize as an \emph{admissibility field}
$A:\cM \to \R_{\geq 0}$ defined over the activation manifold. This object is
not imported as a prior theoretical framework; it emerges from the empirical
observations as the minimal geometric entity capable of simultaneously
explaining all three findings. The central inversion the evidence supports is:
\emph{future structure determines representation}, reversing the standard
assumption that representation determines behavior. We formalize this as the
\emph{Reachability Determines Representation} principle, prove a formal
equivalence relation over activation states, and derive from it a unified field
hypothesis in which activation clusters, ICA directions, effective receptive
fields, and velocity fields in trajectory space all appear as different
measurements of the same underlying object $\fA(x)$. The paper also establishes
a structural correspondence between the three empirical programs and the three
foundational layers of modern differential geometry: topology (which asks what
is connected), differential structure (which asks what directions exist locally),
and dynamics (which asks how trajectories evolve). This triad maps onto Liu's
activation graph, ICALens's independent components, and WMSD's learned velocity
fields almost perfectly, suggesting that the architecture of modern geometry
was, in a precise sense, the right language for activation space all along.
\end{abstract}

\vspace{1.5cm}

{\small
\textbf{Keywords:} activation geometry, admissibility fields, reachability,
mechanistic interpretability, independent component analysis, world models,
representation learning, manifold theory, constraint topology}

\end{titlepage}

\tableofcontents
\newpage

%----------------------------------------------------------------------
\section{Introduction: Three Convergences}
%----------------------------------------------------------------------

Theoretical frameworks in science earn their keep not by fitting a single
observation but by revealing that several apparently unrelated observations
are, in fact, measurements of the same underlying object seen from different
angles. When this happens --- when a minimal formal entity renders transparent
a cluster of empirical findings that had seemed to require separate
explanations --- the result is not merely a new theory but a kind of
perceptual reorganization. What had looked like a landscape of disconnected
facts becomes a topographic map of a single terrain.

This paper argues that something of this kind is occurring right now in the
mechanistic interpretability of large language models. Three independent
empirical programs, developed in near-simultaneity during the first half of
2026, have arrived at structurally identical conclusions by entirely different
routes. The first, Liu's normalization study \cite{liu2026normgeometry},
demonstrates that once the magnitude of residual-stream activations is
discarded via L2 normalization, the resulting points on the unit hypersphere
organize into coherent connected components --- and that those components
correspond not to semantic topics such as ``mammals'' or ``countries'' but to
structural and procedural constraints such as question-and-answer formatting,
copyright boilerplate, patent section headings, LaTeX syntax, and programming
language import blocks. The second, the ICALens study by Liu and Han
\cite{liu2026icalens}, demonstrates that the same activation space, when
subjected to row-normalization followed by whitening and independent component
analysis, reveals directional structure characterized by non-Gaussian,
heavy-tailed projection distributions --- and that these statistically
exceptional directions correspond to modes of textual constraint rather than
to object categories or named concepts. The third, the World Model
Self-Distillation paper by Stapf et al.\ \cite{stapf2026wmsd}, demonstrates
that an instruction-conditioned video world model trained on a rich teacher
signal containing step-by-step solution trajectories can, through a process
of self-distillation combined with reinforcement learning and vision-language
model feedback, learn to match or exceed the teacher's performance using only
high-level task specifications, with additional transfer to robotic tasks --- and that the right mathematical formalization of what the
student is learning is not a set of states or a policy function but a
\emph{velocity field over possibility space}, whose fundamental stability
properties are governed by local matching of continuation dynamics rather than
by state-level imitation.

Each of these results is striking in isolation. Together they are
extraordinary. What they share, as we will argue in careful detail, is a
single structural claim: that the computational geometry of learned
representations is organized primarily around \emph{what can happen next}
rather than around \emph{what is present now}. The activation manifold appears
to be partitioned not into categories of objects but into equivalence classes
of futures --- regions where different tokens, contexts, and computational
states become interchangeable because they license nearly the same
continuation possibilities. The relevant notion of similarity, in other
words, is not lexical or even semantic in the traditional sense; it is
\emph{reachability-theoretic}.

The formal object we introduce to unify these observations is the
\emph{admissibility field} $A:\cM \to \R_{\geq 0}$, where $\cM$ is the
activation manifold. Informally, $A(x)$ measures, at activation state $x$,
how much admissible future trajectory lies ahead --- the measure of the
reachable future set weighted by a characteristic function of valid
continuations. This object is deliberately not introduced as a presupposed
theoretical framework. We introduce it only in Section~4, after establishing
the empirical anomalies in Sections~2 and~3, and we give its full formal
synthesis only in Section~11, after deriving its role in each of the
intermediate sections independently. The strategy throughout is: find the
observations, identify the constraints each places on any explanatory
structure, and determine the minimal formal entity that satisfies all of them.
The admissibility field is our candidate for that entity.

A note on intellectual genealogy. The synthesis we are proposing does not
emerge from nowhere. It draws on a long tradition in manifold learning
\cite{tenenbaum2000isomap, roweis2000lle, coifman2006diffusion}, information
geometry \cite{amari2016information}, percolation theory
\cite{stauffer1994percolation, newman2018networks}, and the emerging field of
mechanistic interpretability \cite{olah2020circuits, elhage2022toy,
bricken2023monosemantic, cunningham2023sae}. It connects to the theoretical
program in world models \cite{ha2018worldmodels, lecun2022path,
stapf2026wmsd} and to the reachability-theoretic ontology developed across
several prior monographs in this series \cite{flyxion2026reachability,
flyxion2026hidden, flyxion2026admissibility, flyxion2026frozen}. The
bibliography at the end of this paper is organized deliberately to reflect the
logical genealogy of the argument: geometry and topology first, interpretability
second, dynamics third, reachability fourth. This ordering is not incidental;
it mirrors the structure of the paper itself and is intended to make clear that
the admissibility field emerges from the empirical literature rather than being
imposed upon it.

The structure of the paper is as follows. Section~2 formalizes Liu's
normalization procedure as a projection operation in the sense of differential
geometry and derives its implications for the structure of angular similarity.
Section~3 examines the empirical clustering results and argues that the content
of the clusters --- procedural and structural constraints rather than semantic
topics --- constitutes a fundamental anomaly that forces a reexamination of
what activation representations encode. Section~4 introduces the Future
Equivalence Principle and the Reachability Determines Representation
inversion, which together provide the theoretical bridge between the clustering
observations and the admissibility field formalism. Section~5 analyzes the
topology of the activation graph, reinterpreting the giant component result and
the power-law component spectrum in terms of the level-set topology of $A(x)$.
Section~6 provides the bridge between graph-theoretic and differential
descriptions, establishing the correspondence between connected components,
ICA directions, and velocity fields as instantiations of the
topology-differential structure-dynamics triad. Section~7 formalizes ICALens's
key technical innovations --- non-Gaussianity as constraint selectivity,
effective receptive field as correlation length, ICA versus SAE --- within the
admissibility geometry framework. Section~8 derives the WMSD velocity-field
formalism and re-reads its key stability theorem as a statement about local
admissibility agreement controlling global reachability. Section~9 formalizes
distillation as an admissibility-preserving projection and establishes the Task
Sufficiency Projection Principle. Section~10 consolidates the ontological
inversion across all three papers, arguing that it is an empirical consequence
rather than a philosophical assertion. Section~11 assembles the Admissibility
Field Hypothesis in its full formal version and identifies the open mathematical
questions the framework raises. Section~12 concludes.

%----------------------------------------------------------------------
\section{Normalization as Projection: The Hyperspherical Turn}
%----------------------------------------------------------------------

\subsection{The Structure of Residual Activations}

Let $h \in \R^d$ be a residual-stream activation of a transformer at some
layer. For the Gemma model analyzed in \cite{liu2026normgeometry}, $d = 2304$.
Each token processed by the model at each layer generates one such vector.
If we collect $n$ tokens, we obtain a point cloud $\{h_1, h_2, \ldots, h_n\}
\subset \R^d$. In Liu's experiment $n \approx 10^6$.

In its raw form, each activation $h_i$ has both a direction and a magnitude:
\begin{equation}
  h_i = \norm{h_i} \cdot \hh_i,
  \qquad
  \hh_i = \frac{h_i}{\norm{h_i}}.
\end{equation}
The full information content of $h_i$ is therefore encoded in the pair
$(\norm{h_i}, \hh_i) \in \R_{>0} \times \Sph^{d-1}$. These two components are
in principle independent, and they carry different kinds of information about
the model's computational state. Magnitude can reflect confidence, salience,
or energy-like properties of the current processing context; direction encodes
the relational and positional structure within the high-dimensional feature
space.

\subsection{Normalization as a Canonical Quotient Map}

L2 normalization discards magnitude and retains only direction. Formally, it
is the map
\begin{equation}
  \pi_{\Sph} : \R^d \setminus \{0\} \lra \Sph^{d-1},
  \qquad
  \pi_{\Sph}(h) = \frac{h}{\norm{h}}.
\end{equation}
This map has a clean algebraic description. The group $\R_{>0}$ of positive
real numbers acts on $\R^d \setminus \{0\}$ by scalar multiplication:
$\lambda \cdot h = \lambda h$ for $\lambda > 0$. Two vectors are in the same
orbit under this action if and only if they have the same direction, i.e.,
$h \sim h'$ iff $h' = \lambda h$ for some $\lambda > 0$. The map
$\pi_{\Sph}$ is precisely the quotient map for this equivalence relation, and
the quotient space is
\begin{equation}
  \bigl(\R^d \setminus \{0\}\bigr) \big/ \R_{>0} \;\cong\; \Sph^{d-1}.
\end{equation}
This algebraic fact has a significant geometric implication: normalization is
not an arbitrary preprocessing step but a geometrically natural operation
corresponding to the identification of scale-equivalent vectors. It reduces the
dimensionality of the relevant structure by one (the scale degree of freedom)
while leaving the directional degrees of freedom intact.

\subsection{Angular Similarity as the Natural Metric}

Once all activations lie on $\Sph^{d-1}$, the natural metric is angular. For
two normalized activations $\hh_i, \hh_j \in \Sph^{d-1}$, the squared
Euclidean distance is
\begin{equation}
  \norm{\hh_i - \hh_j}^2
  = \norm{\hh_i}^2 + \norm{\hh_j}^2 - 2\ip{\hh_i}{\hh_j}
  = 2 - 2\ip{\hh_i}{\hh_j}
  = 2\bigl(1 - \cos(h_i, h_j)\bigr),
\end{equation}
where $\cos(h_i, h_j) = \ip{\hh_i}{\hh_j}$ is the cosine similarity. This
means that Euclidean proximity on the hypersphere is equivalent to cosine
similarity, and that the graph connectivity threshold $\theta$ in Liu's
construction (two activations are connected when their cosine similarity
exceeds $\theta$) corresponds exactly to a geodesic ball condition on the
sphere.

The choice to use angular rather than Euclidean distance is not merely
aesthetic. It reflects a deep assumption about what matters computationally.
The RMSNorm operation used in modern transformers \cite{zhang2019rmsnorm}
normalizes activations before applying learned weight matrices, which means
that the downstream computation at each layer effectively operates on the
normalized direction $\hh_i$ rather than the raw vector $h_i$. The downstream
layers are, in a precise functional sense, insensitive to magnitude and
responsive only to direction. This is why \cite{liu2026normgeometry} argues
that normalization makes activation geometry ``imaginable'': it aligns the
representation space with the space that the computation is actually using.

\subsection{The Astronomical Analogy and Degree-of-Freedom Reduction}

There is a useful analogy from astronomy. When we observe the night sky,
stars appear as points of light that vary enormously in apparent brightness.
Attempting to detect constellations --- geometric patterns of angular positions
--- in the presence of large brightness variation is difficult: a bright nearby
star and a dim distant star may be spatially close on the celestial sphere but
enormously different in apparent brightness. If one were to use raw Euclidean
coordinates in 3-space rather than angular coordinates on the sphere, the
positional structure would be obscured by scale variation.

Normalization performs precisely the same operation as projecting stars onto
the celestial sphere: it discards the radial degree of freedom and reveals the
angular geometry. The constellations become visible not because new information
has been added but because an irrelevant degree of freedom has been removed.
This is the general principle at work. Many representation-learning techniques
\cite{coifman2006diffusion, tenenbaum2000isomap} can be understood as
collapsing irrelevant fibers of a projection in order to expose the geometry
that was present all along.

\subsection{First Structural Observation}

We have established the first structural observation of the paper:
normalization is a geometrically natural projection that collapses the
scale fiber $\R_{>0}$ over each direction, moving activations from
$\R^d \setminus \{0\}$ to $\Sph^{d-1}$ and replacing Euclidean distance with
angular similarity. Crucially, this operation aligns the representation with
the downstream computation, since RMSNorm-based architectures are functionally
insensitive to scale. The surprising empirical result that this projection
\emph{reveals} organized geometric structure rather than destroying it
is the fact that demands explanation. We turn to that structure next.

%----------------------------------------------------------------------
\section{Constraint Clusters, Not Concept Clusters}
%----------------------------------------------------------------------

\subsection{The Graph Construction}

With normalized activations $\{\hh_1, \ldots, \hh_n\} \subset \Sph^{d-1}$
in hand, \cite{liu2026normgeometry} constructs the graph
\begin{equation}
  G_\theta = (V, E_\theta),
  \qquad
  (i,j) \in E_\theta \;\iff\; \ip{\hh_i}{\hh_j} \geq \theta
\end{equation}
for a fixed cosine similarity threshold $\theta$. The vertices $V$ are the
million normalized activations; two vertices are connected by an edge whenever
the angle between the corresponding directions falls below the threshold.
Connected components of $G_\theta$ are then the maximal sets of activations
that are mutually reachable through chains of high-similarity steps.

The central empirical result is the composition of the medium-sized and
high-ranked components. We pause to state this carefully, because the precise
content of the clusters is the observation that drives the entire theoretical
argument.

\subsection{The Central Anomaly}

The na\"{i}ve expectation, grounded in decades of word embedding and language
model analysis, is that activation clusters would correspond to semantic topics.
One expects to find a ``mammals'' cluster, a ``countries'' cluster, an
``emotions'' cluster, a ``mathematical concepts'' cluster. This expectation is
deeply entrenched because it aligns with the distributional hypothesis in
linguistics \cite{rumelhart1986representations} and with the abundant evidence
that language models encode semantic relationships in their representations.

What \cite{liu2026normgeometry} finds is entirely different. The medium-sized
connected components, which are the most informative (the giant component, as
we will analyze in Section~5, is too large to be interpretable as a single
cluster; the isolated singletons are too small), correspond to organizational
and procedural structures of text:

The second-ranked cluster consists entirely of activations corresponding to
the colon character in the pattern ``Q:'', i.e., the boundary between a
question header and its content in question-and-answer formatted text. The
other significant clusters correspond to copyright notice headers, patent
section headings of the form ``Field of the Invention'' and similar, legal
citation structures, software license boilerplate of the form ``GNU General
Public License'', LaTeX mathematical notation delimiters and structural
commands, JSON escape sequences and structural punctuation, Go language
indentation patterns, and Java import block headers.

Let us be precise about why this is anomalous. These are not
semantic topics. They share no common subject matter. A patent heading has
nothing in common with a JSON escape sequence in terms of what it \emph{refers
to}. A copyright notice and a LaTeX fragment do not co-occur in documents, do
not describe related concepts, and do not involve similar vocabulary. If
semantic similarity were driving the clustering, we would expect such items to
appear in entirely different parts of the activation space.

What they share is something else entirely: they are all located at boundaries
into \emph{constrained continuation regimes}. Each of these textual structures
is a gateway into a region of text space where the next tokens are highly
predictable and the range of admissible continuations is severely restricted.

\subsection{The Rank-2 Result in Detail}

The second-ranked cluster --- 276 activations all corresponding to the colon
after ``Q'' --- deserves especially careful analysis. The colon ``\texttt{:}''
is an extremely common character that occurs in thousands of different textual
contexts: after labels in YAML files, in URL structures, in time expressions,
after Greek citations, after ratios, in emoji, in mathematical notation, and
in ordinary English prose. The fact that cosine similarity in normalized
activation space selects out precisely and only the subset of colons that appear
in question-and-answer context is not trivially explained by either the
token's lexical identity or its syntactic role.

The 276 instances of ``Q:'' in the dataset come from many different documents,
many different topics, many different authors, many different styles. They share
precisely this: after ``Q:'', the model is in a highly specific
computational state. The continuation is constrained to be a question or
question-like text. The format that follows is expected to have specific
structural properties. The future trajectory of the text, in a precise
mathematical sense, is more constrained than it is for most other contexts.
The activation for the colon after ``Q'' is not encoding the identity of the
colon. It is encoding a location in \emph{future-space}: the state of being
at the entrance to a constrained continuation manifold.

\begin{observation}
  The Rank-2 cluster in Liu's activation graph demonstrates that the network
  has learned activation states that are not attached to token identity but to
  the structure of the future continuation they initiate. The token ``\texttt{:}''
  merely serves as an entrance into a region of state space; the activation
  encodes the properties of that region, not the properties of the token.
\end{observation}

This observation is not limited to the Rank-2 cluster. It applies uniformly
to the identified components. The activation corresponding to ``Field of the
Invention'' is useful to the model not because it must identify patent headings
as a lexical category but because, from that activation state, the admissible
continuations are a very particular set: technical descriptions of inventions,
problem statements, prior art summaries, legal claims, and similar
patent-register prose. The activation corresponding to ``import ('' in Go
source code is at a boundary into a structured block where the admissible
tokens are package identifiers and structural delimiters. The activation
corresponding to copyright boilerplate initiates a structured sequence whose
continuation is almost entirely predictable.

\subsection{Process Attractors versus Noun Clusters}

The distinction being drawn here has a natural formulation in dynamical systems
language. Classical attractor analysis in dynamical systems identifies regions
of state space toward which trajectories converge under the governing dynamics.
We can analogize: the activation clusters discovered by \cite{liu2026normgeometry}
are regions of activation space that are not distinguished by the object they
represent but by the dynamic they initiate. They are process attractors rather
than noun clusters.

A noun cluster would be a region of activation space that all activations for
``dog''-related tokens converge to because dogs share common features. A
process attractor is a region that activations converge to whenever the text
enters a particular mode of structured continuation, regardless of the
subject matter being discussed. The Q/A cluster is a process attractor for the
beginning-of-question state; the patent-heading cluster is a process attractor
for the beginning-of-patent-section state; the copyright cluster is a process
attractor for the beginning-of-boilerplate state.

\subsection{Second Structural Observation}

This gives us the second structural observation: the geometric structure
revealed by normalization is organized around continuation constraints rather
than denotational categories. Activations cluster together not because they
refer to similar things but because they initiate similar kinds of future
trajectories. This is our central empirical anomaly, and it demands a
theoretical account. In the next section we provide one.

%----------------------------------------------------------------------
\section{The Future Equivalence Principle and the Admissibility Field}
%----------------------------------------------------------------------

\subsection{Why Cosine Similarity Is the Wrong Metric}

The graph $G_\theta$ uses cosine similarity as its proximity criterion.
This is a natural choice given that the activations have been projected onto
the hypersphere, and it is the criterion that reveals the clustering structure
described in Section~3. But the empirical content of the clusters suggests that
cosine similarity is not the fundamental similarity relation at work; it is,
rather, a proxy for a deeper relation.

To see this, consider the following thought experiment. Suppose we have two
activations $a_1$ and $a_2$ that are far apart in angular distance on
$\Sph^{d-1}$ but that both occur at the entrance to patent-heading continuation
regimes. Standard cosine similarity would classify them as dissimilar. But from
the perspective of what the model will do next --- what futures are available ---
they are functionally identical. Conversely, suppose $a_1$ and $a_2$ are
cosine-similar in activation space but one occurs at a patent heading and one
occurs in a general prose context. Cosine similarity would classify them as
similar, but they license entirely different continuation manifolds.

The correct similarity relation is not determined by the current geometry of
the activation vectors. It is determined by the geometry of the futures they
generate. This motivates the central formal construction of the paper.

\subsection{The Reachability Operator}

\begin{definition}[Reachability Operator]
  Let $a \in \cM$ be an activation state, where $\cM \subset \Sph^{d-1}$ is the
  normalized activation manifold. Let $\tau = (a_{t_0}, a_{t_1}, \ldots,
  a_{t_k})$ denote a future activation trajectory beginning at $a_{t_0} = a$.
  The \emph{reachability operator} at $a$ is defined as
  \begin{equation}
    \reachset{a} = \cL(\tau \mid a),
  \end{equation}
  the probability distribution over future activation trajectories $\tau$
  conditioned on the current state $a$. Here $\cL(\cdot \mid a)$ denotes the
  law of the trajectory process induced by the model's continuation dynamics
  starting from $a$.
\end{definition}

Informally, $\reachset{a}$ is the distribution over futures accessible from
$a$. It encodes not merely the single most likely continuation but the entire
manifold of continuation possibilities, weighted by their probability under
the model's learned dynamics.

\subsection{Future Proximity and Exact Equivalence}

\begin{definition}[$\varepsilon$-Future Proximity]
  Let $D(\cdot, \cdot)$ be a distance on the space of trajectory distributions
  (for concreteness, one may take $D$ to be the Wasserstein-2 distance or the
  KL divergence). For $\varepsilon \geq 0$, define the
  \emph{$\varepsilon$-future proximity relation} on $\cM$ by
  \begin{equation}
    a_1 \approx_R^\varepsilon a_2
    \quad\iff\quad
    D\bigl(\reachset{a_1},\, \reachset{a_2}\bigr) \leq \varepsilon.
  \end{equation}
  In the limit $\varepsilon \to 0$ this recovers exact future equivalence:
  $a_1 \simR a_2 \iff \reachset{a_1} = \reachset{a_2}$.
\end{definition}

\begin{remark}[Transitivity]
  For $\varepsilon > 0$, the relation $\approx_R^\varepsilon$ is reflexive and
  symmetric but generally \emph{not transitive}: one can have
  $a_1 \approx_R^\varepsilon a_2$ and $a_2 \approx_R^\varepsilon a_3$ without
  $a_1 \approx_R^\varepsilon a_3$, since the two balls of radius $\varepsilon$
  may overlap at $a_2$ while being separated at $a_1$ and $a_3$. The relation
  is therefore a \emph{uniform tolerance structure} (a proximity relation in the
  sense of Zeeman) rather than a classical equivalence relation. Genuine
  equivalence classes partition $\cM$ only in the exact limit $\varepsilon = 0$,
  yielding the quotient $\cM / {\simR}$ whose points are maximal future
  equivalents. Throughout the paper, ``equivalence class'' refers to this
  exact-limit partition; for finite $\varepsilon$ we speak of \emph{proximity
  neighborhoods}.
\end{remark}

\begin{principleBox}
\textbf{Future Equivalence Principle.}
Two activation states are functionally equivalent to the extent that they
preserve the same set of admissible continuations. The natural similarity
structure on activation states is future proximity $\approx_R^\varepsilon$,
which recovers a genuine equivalence relation $\simR$ only in the exact limit
$\varepsilon \to 0$. Cosine similarity is a proxy for this deeper structure,
not the fundamental metric.
\end{principleBox}

The motivation for this principle is the following informal but compelling
argument. A transformer activation $a$ is useful to the model exclusively
insofar as it contributes to the computation of subsequent activations and
ultimately to the probability distribution over next tokens. If two activations
$a_1$ and $a_2$ induce nearly identical distributions over future activation
trajectories --- if $D(\reachset{a_1}, \reachset{a_2}) \leq \varepsilon$ for
small $\varepsilon$ --- then no downstream observer restricted to observing the
model's behavior can distinguish between them: they produce the same
probability distributions over any observable output. They are, in the most
operationally meaningful sense, the same state.

Conversely, if $D(\reachset{a_1}, \reachset{a_2})$ is large, then $a_1$ and
$a_2$ generate observably different futures, even if they happen to be close
in activation space by cosine similarity. Two states that are geometrically
proximate but functionally distant --- like a colon in prose versus a colon in
``Q:'' --- represent different computational situations despite their
superficial angular similarity.

\subsection{The Reachability Determines Representation Conjecture}

The Future Equivalence Principle gives us a refined notion of similarity. But
it also opens the door to a much stronger claim. The empirical evidence from
Section~3 suggests not merely that reachability \emph{should} be used as a
similarity criterion, but that the network has \emph{learned} to organize its
representation space according to future proximity classes. That is, the
network's geometry has been shaped by its own future-prediction task in such
a way that future-proximate states end up geometrically proximate.

An important clarification is required before stating this claim. The evidence
in Section~3 establishes primarily that geometrically clustered states share
similar continuation structures: clustering $\Rightarrow$ similar futures.
The claim in the other direction --- similar futures $\Rightarrow$ geometric
clustering --- is the stronger, causal claim that the following conjecture
asserts. The evidence motivates the conjecture but does not yet establish it.

\begin{principleBox}
\textbf{Reachability Determines Representation (RDR Conjecture).}
If $\reachset{a_1} \approx \reachset{a_2}$, then $a_1$ and $a_2$ tend to
become geometrically close in the learned activation space. Formally:
\begin{equation}
  D\bigl(\reachset{a_1}, \reachset{a_2}\bigr) \leq \varepsilon
  \quad\Ra\quad
  \norm{\hh_{a_1} - \hh_{a_2}} \leq f(\varepsilon),
\end{equation}
where $f$ is a monotone increasing function with $f(0) = 0$ that reflects the
degree to which training geometry tracks reachability geometry.

\emph{Falsification criterion.} The RDR conjecture is falsified if two
activation states with nearly identical continuation distributions
($D(\reachset{a_1}, \reachset{a_2}) \leq \varepsilon$ for small $\varepsilon$)
are found to be systematically distant in the activation geometry. The
existence of such pairs, if common, would mean that reachability similarity
is necessary but not sufficient for geometric proximity, and some other factor
determines the representational distance.
\end{principleBox}

This conjecture inverts the standard assumed causal direction in representation
learning. The traditional view is:
\begin{equation}
  \text{Representation} \;\lra\; \text{Behavior}.
\end{equation}
The token has a feature vector; that feature vector, through the transformer's
computations, determines the probability distribution over next tokens.
Representation is prior; behavior is derived.

The RDR conjecture reverses this:
\begin{equation}
  \text{Future Structure} \;\lra\; \text{Representation}.
\end{equation}
The relevant structure of the future continuation manifold --- what is
reachable, what is constrained, what is opened or closed by the current state ---
determines where in activation space the state ends up. Representation is
derived; future structure is prior. The evidence from Section~3 is consistent
with this conjecture; establishing it in either direction would require
constructing pairs of states and comparing their reachability distance with
their activation-space distance, a comparison the framework now makes
precisely actionable.

\subsection{The Admissibility Field: First Introduction and Its Limitations}

We are now in a position to introduce the central formal object of the paper.
The reachability operator $\reachset{a}$ is a rich object --- a full
probability distribution over trajectory space --- and is difficult to work
with directly. The natural first step is to summarize it by a scalar field.
However, before doing so we must acknowledge an important subtlety.

The empirical observations described in Section~3 organize around
\emph{constrained continuation regimes}: copyright notices, patent headings,
Q/A separators, import blocks. These are not characterized primarily by
\emph{validity} (all are perfectly valid text) but by \emph{constraint density}
--- the degree to which the reachable future space is narrow, specific, and
sharply bounded. A copyright notice and a generic prose sentence are equally
``valid,'' yet the activation geometry treats them as categorically different
because their continuation manifolds have radically different shapes.

This suggests that the natural scalar summary of $\reachset{a}$ is not the
probability of validity but something closer to the \emph{geometry} of the
reachable future set. Two natural candidates are the volume of the reachable
set and its entropy:
\begin{equation}
  \fA_{\mathrm{vol}}(x) = \operatorname{Vol}(R(x)),
  \qquad
  \fA_{\mathrm{ent}}(x) = H\bigl(\reachset{x}\bigr),
\end{equation}
where $\operatorname{Vol}$ is an appropriate measure on trajectory space and
$H$ is the differential entropy of the future distribution. Highly constrained
states (patent headings, copyright notices) would have small $\fA_{\mathrm{vol}}$
and small $\fA_{\mathrm{ent}}$; unconstrained generic states would have large
values of both. These summaries are arguably closer to what the activation
geometry is actually measuring than a probability of validity.

More fundamentally, the richer object underlying all of these scalar summaries
is the \emph{reachability bundle}:
\begin{equation}
  \mathcal{R} = \bigl\{ R(x) \bigr\}_{x \in \cM},
\end{equation}
the collection of all reachability distributions, indexed by points on the
activation manifold. The scalar admissibility field $\fA$ is any scalar
functional of $\mathcal{R}$ --- one observable derived from a higher-order
geometric object. The reachability bundle is the fundamental structure; the
admissibility field is a projection. This observation, which we develop further
in Section~11, suggests that the paper's central object may itself be a derived
quantity of a deeper geometry.

For the purposes of organizing the present synthesis, we work with a general
scalar admissibility field $\fA : \cM \to \R_{\geq 0}$ without committing to
a specific choice of functional form. The formal definition used throughout
is the measure of the admissible trajectory set from each state, with the
understanding that this can be refined to volume, entropy, or other geometric
quantities in future work.

\begin{definition}[Admissibility Field]
  Let $\cF_{\mathrm{valid}}$ be a measurable set of continuation trajectories
  considered admissible in a given context. The \emph{admissibility field} is
  the scalar function
  \begin{equation}
    \fA : \cM \lra \R_{\geq 0},
    \qquad
    \fA(x) = \mu_x\bigl(\reachset{x} \cap \cF_{\mathrm{valid}}\bigr),
  \end{equation}
  where $\mu_x$ is the probability measure on trajectory space induced by the
  model's continuation dynamics starting from $x$. The choice of
  $\cF_{\mathrm{valid}}$ is context-dependent; empirically, the activation
  geometry suggests that the relevant criterion is structural constraint density
  (the shape of the future space) rather than semantic validity in any narrow
  sense.
\end{definition}

This field will be the organizing object of the rest of the paper. We preview
its roles: activation clusters will appear as level sets or basins of $\fA$
(Section~5), ICA directions will appear as preferred gradient directions of
$\fA$ (Sections~6 and~7), and WMSD's velocity fields will appear as flows
through $\fA$ (Sections~8 and~9). The full synthesis, including the reachability
bundle perspective, will be assembled in Section~11.

%----------------------------------------------------------------------
\section{The Giant Component and Admissibility Topology}
%----------------------------------------------------------------------

\subsection{The Component Spectrum}

The full component spectrum of Liu's graph $G_\theta$ contains roughly $10^6$
activations. The dominant feature is a single giant connected component
containing approximately 826,000 activations --- about 82.6\% of the total.
The remaining 17.4\% of activations are distributed across a large number of
smaller components, ranging from mid-sized clusters of hundreds to thousands of
activations down to isolated singletons.

A first reaction might be that the giant component is uninteresting --- that
it is simply the ``general language'' blob and the interesting structure is
entirely in the smaller components. This reaction is understandable but misses
something important. The structure of the \emph{entire} spectrum, including the
relative sizes of the giant component and the distribution of the smaller
components, is itself informative. After removing the giant component, the
residual component size distribution appears approximately power-law:
\begin{equation}
  P(\text{component size} = k) \sim k^{-\alpha}
\end{equation}
for some exponent $\alpha > 1$. Power-law distributions in graph component
spectra are a signature of critical phenomena.

\subsection{The Percolation Analogy}

The emergence of a giant connected component in a network with a power-law
residual spectrum is the hallmark of a percolation transition. In percolation
theory \cite{stauffer1994percolation}, a random geometric graph on $n$ points
in $\R^d$ --- where two points are connected whenever they are within distance
$r$ of each other --- undergoes a phase transition as $r$ varies. Below the
critical threshold $r_c$, the graph consists of many small disconnected
clusters. Above $r_c$, a single giant connected component emerges that
contains a finite fraction of all vertices in the infinite-volume limit. At
$r_c$ itself, the component size distribution follows a power law.

The structure Liu observes is consistent with the graph being near or at such
a critical threshold. The presence of a giant component containing the
overwhelming majority of activations, combined with a power-law tail of small
specialized components, suggests that the threshold $\theta$ is set in the
vicinity of the percolation critical point for this particular activation
distribution. This has a significant physical interpretation.

\subsection{Admissibility Density and the Topological Decomposition}

We can now interpret the component spectrum in terms of the admissibility field
$\fA$. Define an \emph{admissibility density} on $\cM$ as follows: high-$\fA$
regions are those where the continuation constraints are sufficiently universal
that activations from many different contexts have overlapping reachability
operators and thus cluster together. Low-$\fA$ regions are those where the
continuation is so highly constrained and specific that only activations from
very similar specialized contexts have overlapping reachability operators.

The topological decomposition of the activation manifold is then:
\begin{equation}
  \cM = G \cup \bigcup_{i=1}^\infty \cI_i,
\end{equation}
where $G$ is the giant component (the generic semantic ocean) and the $\cI_i$
are the specialized admissibility islands. The giant component consists of
states where $\fA$ is smoothly varying and broadly connected: generic prose,
general language, unrestricted continuation contexts. The islands consist of
states where $\fA$ changes rapidly over short distances because the continuation
constraints shift abruptly from those of the generic language into a specialized
mode.

This gives a precise relationship between the component structure and the
admissibility field:
\begin{equation}
  \text{component structure of } G_\theta
  \;\approx\;
  \text{level-set topology of } \fA(x).
\end{equation}
The level sets $\{x : \fA(x) = c\}$ for different values of $c$ trace out the
boundaries between regions of different constraint density. The connected
components of $G_\theta$ are approximations to the path-connected components
of the superlevel sets $\{x : \fA(x) \geq c\}$ for an appropriate $c$.

\subsection{Admissibility Islands as Linguistic Ecosystems}

The specialized islands identified in Liu's analysis --- patent language, legal
citation structures, software licenses, LaTeX equation environments ---
are revealing. Each constitutes what we might call a \emph{linguistic ecosystem}:
a sub-register of text with unusually strong internal constraints and unusually
low coupling to the surrounding activation space. The internal constraint density
is high: given a state inside the patent-language island, the next state is
almost certainly also inside that island. The external coupling is low: it is
difficult to reach the patent-language island from generic prose contexts or
from programming language contexts without passing through specific boundary
tokens.

These properties define an isolated submanifold in activation space. The
activations corresponding to patent text form a cluster not because patent
text is semantically homogeneous --- it is not, covering engineering, chemistry,
biology, software, and countless other subjects --- but because the
\emph{structural rules} governing patent text create a uniform admissibility
landscape across all those subject matters. The level sets of $\fA$ cut across
subject-matter distinctions and group together anything that sits inside the
same constraint field.

\subsection{Connection to RSVP Phase Structure}

The percolation picture also connects naturally to phase structure in field
theories. In the RSVP framework \cite{flyxion2026admissibility}, the
admissibility field can undergo phase transitions as the parameter governing
constraint density crosses a critical value. Near the critical point, the
correlation length diverges and the component spectrum becomes scale-free ---
which is precisely the power-law behavior observed. The giant component is the
``high-density phase'' of the admissibility field; the isolated islands are
``nucleated phases'' corresponding to local maxima of the constraint potential
embedded in a lower-density background.

\subsection{Third Structural Observation}

The third structural observation is: the component spectrum of the activation
graph encodes the level-set topology of the admissibility field $\fA$. The
giant component corresponds to the high-connectivity basin of $\fA$; the
specialized islands correspond to isolated high-constraint regions where the
admissibility landscape is steep and the crossing energy is high. The
power-law distribution of component sizes is consistent with the graph being
near a percolation critical point.

%----------------------------------------------------------------------
\section{Constraint Fields and Natural Coordinates: A Bridge}
%----------------------------------------------------------------------

\subsection{Two Descriptions of the Same Geometry}

At this point in the argument we have two empirical programs --- Liu's
activation graph and Liu \& Han's ICA decomposition --- and a single theoretical
object $\fA$. It might appear that the graph-based and ICA-based descriptions
are independent discoveries that require separate theoretical accounts. In this
section we argue that they are in fact dual descriptions of the same geometric
object, seen from two fundamentally different perspectives: one topological and
one differential.

The key mathematical distinction is the following. A connected component of a
graph is a \emph{topological} object: it depends only on the qualitative
question of whether a path exists between two vertices, not on the quantitative
local geometry. A direction recovered by ICA is a \emph{differential} object:
it is a tangent direction to the manifold at a point, and its relevance depends
on how the distribution varies locally. The graph asks which points belong
together; ICA asks along which directions the distribution fractures.

\subsection{Level Sets and Gradient Directions}

Suppose the activation manifold $\cM$ is embedded in $\Sph^{d-1}$ and the
admissibility field $\fA : \cM \to \R_{\geq 0}$ is a smooth scalar function
on this manifold. The gradient $\nabla \fA(x)$ at a point $x \in \cM$ is a
tangent vector pointing in the direction of steepest ascent of $\fA$. The
level sets $\{x : \fA(x) = c\}$ are the $(d-2)$-dimensional hypersurfaces
perpendicular to $\nabla \fA$.

Connected components of the graph $G_\theta$ are approximating the connected
components of the superlevel sets of $\fA$. These are precisely the level-set
topology. ICA directions, on the other hand, are seeking the directions in
the tangent space along which the projection distribution is maximally
non-Gaussian. We will argue in Section~7 that constraint selectivity implies
heavy-tailed projections, which means that ICA directions are approximating
the directions of largest variation in $\fA$:
\begin{equation}
  u_k \;\approx\; \text{directions of large variation in } \nabla \fA(x).
\end{equation}
Graph components thus reveal the topology of $\fA$ --- which regions are
connected, which are separated --- while ICA directions reveal the differential
structure of $\fA$ --- where it varies most rapidly and in which directions.
They are complementary instruments for probing the same geometric object.

\subsection{The Fundamental Geometric Triad}

This duality positions us to state the central structural observation of the
bridge section. The three empirical programs of this paper instantiate exactly
the three foundational layers of modern differential geometry:

\emph{Topology} asks what is connected, what is separated, what paths exist
between points, how many holes a space has. This is the domain of Liu's
activation graph: which activations belong to the same component, which
components are isolated, what is the genus of the resulting graph structure.

\emph{Differential geometry} asks what directions exist locally, what the
curvature is, how the tangent spaces at different points are related. This is
the domain of ICALens: which directions are preferred by the distribution,
what is the statistical geometry of the local tangent space, which directions
are most selective.

\emph{Dynamics} asks how trajectories evolve through the space, what the vector
fields are, what the attractors and repellers are. This is the domain of WMSD:
what velocity field governs motion through possibility space, how that field is
learned and transferred, what its stability properties are.

The correspondence is:
\begin{equation}
  \text{Topology} \;\lra\; \text{Differential Structure} \;\lra\; \text{Dynamics}
\end{equation}
\begin{equation}
  \text{Liu's graph} \;\lra\; \text{ICALens directions} \;\lra\; \text{WMSD velocity fields}
\end{equation}
and all three have the admissibility field $\fA$ as their common underlying
object:
\begin{equation}
  \text{graph components} \;\approx\; \text{level sets of } \fA, \qquad
  \text{ICA directions} \;\approx\; \text{gradient directions of } \fA, \qquad
  \text{velocity fields} \;\approx\; \text{flows through } \fA.
\end{equation}

This triad is not merely an aesthetic observation. Modern differential geometry
is organized precisely in this order --- topology first, differential structure
second, dynamics third --- because each layer presupposes and refines the
previous one. Topology establishes what regions exist; differential geometry
adds local quantitative structure to those regions; dynamics adds the temporal
evolution that connects points within and between regions. The fact that three
independent empirical programs in mechanistic interpretability and world modeling
map onto these three layers is strong circumstantial evidence that the
mathematical structure of the activation manifold is genuinely geometric in a
deep sense.

%----------------------------------------------------------------------
\section{ICA as Natural Cleavage: Statistical Fractures in Activation Space}
%----------------------------------------------------------------------

\subsection{The ICALens Setup}

The ICALens paper \cite{liu2026icalens} takes a different entry point into the
same geometry. Rather than constructing a graph over activations, it applies
Independent Component Analysis (ICA) \cite{hyvarinen2000ica, hyvarinen2001book}
to the normalized and whitened activation matrix. The paper evaluates this
approach across three model families: GPT-2 Small ($d = 768$), Gemma~2~2B
($d = 2304$), and Qwen~3.5~2B~Base ($d = 2048$), fitting ICA independently at
each residual-stream layer using one million token positions drawn from the
Pile-10k dataset.

The preprocessing pipeline addresses a fundamental challenge: raw transformer
activations are highly anisotropic, containing systematic outlier dimensions,
rare massive activations, and attention-sink tokens that destabilize standard
ICA fitting. The ICALens pipeline proceeds in three steps. First, each
activation vector is row-normalized by its $\ell_2$ norm:
\begin{equation}
  r(x_i) = \frac{x_i}{\max(\norm{x_i}_2, \varepsilon)},
\end{equation}
reducing the influence of activation-norm outliers before any further
processing. Second, the row-normalized matrix is centered and whitened to
remove second-order correlations. Third, a GPU-parallel FastICA algorithm is
applied to the whitened matrix with a robust convergence criterion (the
p95-LIM rule, which accepts a fit when 95\% of components have stabilized
even if a small tail remains difficult to converge) and an adaptive refit
procedure that reduces the target component count for layers where full-rank
convergence fails.

This engineering makes the difference: on GPT-2 Small with one million
activations, the ICALens pipeline increases the number of accepted layers by
400\% compared to a na\"{i}ve FastICA implementation, and the resulting
components are validated through human annotation, targeted prompt tests,
and SAEBench evaluations.

The key interpretive claim of the paper is: \emph{interpretable directions are
statistically exceptional directions}, and statistical exceptionality is
measured by non-Gaussianity of the scalar projection distribution. The paper
also establishes empirically that SAE decoder directions, despite being trained
for sparse reconstruction rather than non-Gaussianity, already exhibit elevated
kurtosis relative to random directions --- suggesting that non-Gaussianity is a
common statistical signature of learned feature directions, which ICA makes
explicit as its direct optimization target.

\subsection{Non-Gaussianity as Constraint Selectivity}

Let $v \in \Sph^{d-1}$ be a candidate direction and let $z_i(v) = \ip{v}{\hh_i}$
be the scalar projection of the $i$-th normalized activation onto $v$. Define
the excess kurtosis of this projection as
\begin{equation}
  \kur(v) =
  \frac{\frac{1}{n}\sum_{i=1}^n \bigl(z_i(v) - \bar{z}(v)\bigr)^4}
  {\left(\frac{1}{n}\sum_{i=1}^n \bigl(z_i(v) - \bar{z}(v)\bigr)^2\right)^2}
  - 3,
\end{equation}
where $\bar{z}(v) = \frac{1}{n}\sum_i z_i(v)$ is the mean projection. For a
Gaussian distribution, $\kur(v) = 0$; heavy-tailed distributions have
$\kur(v) > 0$. FastICA maximizes $|\kur(v)|$ (or a smooth approximation thereof)
to find the most non-Gaussian projection directions.

Why should interpretable directions be non-Gaussian? The argument runs as
follows. If a direction $v$ is relevant to a specific constraint regime ---
say, the ICA direction for question-and-answer boundaries --- then for the vast
majority of tokens in the dataset, that direction will project to values near
zero (since most tokens are not at Q/A boundaries). For the small fraction of
tokens that are at Q/A boundaries, the projection will be large. The resulting
distribution is quiet for most tokens and extreme for a few: precisely a
heavy-tailed, high-kurtosis distribution.

By contrast, a direction that is relevant to generic semantic content would
be activated moderately across a large fraction of tokens --- different tokens
activate it to different degrees, but not with the extreme sparsity of
constraint-regime directions. Such a direction would have a projection
distribution closer to Gaussian.

The chain of implication is therefore:
\begin{equation}
  \text{constraint-regime selectivity}
  \;\Ra\;
  \text{sparse activation}
  \;\Ra\;
  \text{heavy-tailed projection}
  \;\Ra\;
  \text{ICA recoverability}.
\end{equation}
ICA recovers constraint-regime-selective directions not because it was
designed with constraints in mind, but because the statistical footprint of
constraint selectivity --- sparsity --- produces the non-Gaussianity that ICA
maximizes.

\subsection{ICA Directions as Gradient Directions of the Admissibility Field}

We can now connect the ICA directions to the differential structure of $\fA$.
In the differential-geometric picture, the admissibility field $\fA$ varies
smoothly over the activation manifold $\cM$. At any point $x \in \cM$, the
gradient $\nabla \fA(x)$ points in the direction of fastest increase of $\fA$.
Directions along which $\fA$ changes rapidly are directions where the
constraint landscape is unstable: crossing a small distance in that direction
moves from one constraint regime to a qualitatively different one.

ICA's non-Gaussianity criterion selects directions that are most selective
between high-$\fA$ and low-$\fA$ states. For a direction $u$ perpendicular to
a level set of $\fA$, the projection $z_i(u)$ will be large for activations
on one side of the level set and small for activations on the other side: this
is precisely the sparse, bimodal structure that produces high kurtosis.

Therefore, ICA directions approximate the directions of large variation in
$\nabla \fA$ at representative points:
\begin{equation}
  u_k \;\approx\; \underset{v \in \Sph^{d-1}}{\arg\max}
  \; \mathbb{E}_{x \in \cM}\bigl[\norm{\nabla_v \fA(x)}^2\bigr],
\end{equation}
where $\nabla_v \fA(x) = \ip{\nabla \fA(x)}{v}$ is the directional derivative
along $v$. The ICA decomposition is, in this light, approximating the
eigendirections of the Hessian of the expected admissibility variation --- the
natural coordinate axes of the admissibility landscape.

\subsection{Effective Receptive Field as Correlation Length}

A significant technical contribution of \cite{liu2026icalens} is the
\emph{effective receptive field} (ERF) diagnostic. ERF asks a precise
question: how much left context is sufficient to recover the same signed
component response at a target token? For a component $j$, an evidence set
$\mathcal{E}_j$ is formed from examples whose absolute component score
exceeds half the maximum score for that component. For each evidence example
$(x, t)$ in $\mathcal{E}_j$, suffixes of increasing length
$x_{t-k+1:t}$ for $k = 1, 2, \ldots, K_{\max}$ are constructed and the
component is recomputed. The sample-level ERF is the shortest suffix length
$k$ such that component $j$ remains among the top-15 most active components
and preserves its original sign:
\begin{equation}
  \mathrm{erf}_j(x, t) = \min\{k \in \{1, \ldots, K_{\max}\} :
  j \in \mathrm{Top}_{15}\!\left(\bigl\{|s_r(h_t^{(k)})|\bigr\}_r\right)
  \;\wedge\;
  \mathrm{sign}(s_j(h_t^{(k)})) = \mathrm{sign}(s_j(h_t(x_{1:t})))\}.
\end{equation}
The component-level ERF is the mean over the evidence set:
\begin{equation}
  \ERF(j) = \frac{1}{|\mathcal{E}_j|}
  \sum_{(x,t) \in \mathcal{E}_j} \mathrm{erf}_j(x, t).
\end{equation}
The paper uses $K_{\max} = 11$ and evaluates across all three model families.
Components with small ERF are sensitive only to the current token or a short
local window; components with large ERF depend on broader discourse context.

The paper reports a layer-wise shift across all three model families:
components in shallow layers are dominated by token-local directions
(surface constraints --- tokenization, local syntax, character-level
patterns), while components in deeper layers show more medium- and
long-context directions. Importantly, the paper notes this is a gradual
shift in mixture proportions rather than a sharp handoff, and that the
largest-ERF components (those requiring more than the 11-token window)
are most prevalent in \emph{middle} layers --- making middle layers a
particularly useful target for inspecting context-dependent structure.
The paper also finds a consistent negative correlation between ERF and
excess kurtosis (Spearman correlations of $-0.41$ to $-0.50$ across
models): high-kurtosis components tend to have small ERF, because they
are activated by narrow recurring patterns whose triggering condition is
visible at the target token, making them easier to annotate and explain.

This is precisely the local-to-global transition predicted by the
topology-differential structure-dynamics triad. Shallow layers are where the
topology of the activation manifold is established --- which tokens are
connected to which. Deeper layers are where the dynamics play out --- where
long-range constraint fields become activated and govern the trajectory of the
text over extended sequences.

In field-theoretic language, ERF is a measure of the \emph{correlation length}
$\xi_k$ of the $k$-th ICA component. The correlation length of a field
measures the spatial (or here, contextual) range over which fluctuations in
that field are correlated. A component with $\ERF(u_k) = 5$ tokens has
$\xi_k \approx 5$: its activation is determined by local context and decays
rapidly with context length. A component with $\ERF(u_k) = 100$ tokens has
$\xi_k \approx 100$: its activation is shaped by the document structure over
long ranges.

The layer-wise shift from small to large ERF then corresponds to the
well-known renormalization-group picture of a field theory: short-range
fluctuations are integrated out in early layers, leaving only the long-range
modes that govern macro-scale behavior. This connection deserves more formal
development than we can provide here; we flag it as one of the open questions
in Section~11.

\subsection{ICA versus SAE: Natural Coordinates versus Trained Dictionaries}

The \cite{liu2026icalens} paper positions ICA as a ``first lens'' before
training sparse autoencoders (SAEs) \cite{cunningham2023sae,
bricken2023monosemantic}. The comparison is instructive for our purposes.

A sparse autoencoder learns a dictionary $\{d_1, \ldots, d_m\} \subset \R^d$
with $m \gg d$ such that each activation can be expressed as a sparse
combination: $h \approx \sum_{k \in S} c_k d_k$ for a small index set $S$.
The dictionary is trained to optimize a reconstruction-plus-sparsity objective,
and the resulting features are interpretable as learned semantic or functional
categories. Crucially, the dictionary is an artifact of the training procedure;
the directions $d_k$ have no privileged status as natural axes of the
activation manifold independent of the SAE training.

ICA, by contrast, recovers directions that are statistically privileged in the
\emph{original} activation distribution --- directions along which the marginal
distribution of projections is maximally non-Gaussian. These directions are
not learned by an auxiliary model; they are properties of the activation
distribution itself. In the language of differential geometry, ICA is finding
the natural cleavage planes --- the directions along which the manifold most
naturally fractures --- while SAEs are constructing a learned coordinate system
that may or may not align with those natural directions.

The empirical comparison reported in \cite{liu2026icalens} is more nuanced
and more favorable to ICA than one might expect. On SAEBench sparse probing
across the eight standard concept datasets, ICA is competitive with
publicly released SAE dictionaries (Gemma Scope, Qwen Scope) and
consistently outperforms both ITDA (a training-light sparse-coding
alternative) and PCA. Compact Matryoshka SAE variants with smaller
dictionary sizes also remain below ICA under equivalent component budgets,
which means the ICA advantage is not simply a function of dictionary size.
On SAEBench Targeted Probe Perturbation --- an intervention-based
disentanglement metric --- ICA is strongest relative to public SAEs at
small-to-medium intervention budgets, where a small number of ICA components
captures enough class-relevant information to support selective interventions
without requiring search through a large overcomplete dictionary.

The directional overlap analysis also reveals that ICA and SAE are
complementary rather than redundant. Most ICA components have a nearest-SAE
neighbor with moderate cosine similarity (roughly 0.3--0.6 in the median
across layers and models) rather than near-perfect alignment, indicating
partial but not complete agreement. Moreover, SAE features often activate
on localized tokens (consistent with the sparsity pressure in SAE training),
while ICA components frequently vary more smoothly across spans of related
tokens --- tracking, for example, the financial interpretation of ``bank''
across deposit, check, withdraw, and cash in the same sentence. This
difference in activation pattern is not accidental: it reflects the different
objectives of the two methods. SAEs maximize sparse reconstruction; ICA
maximizes non-Gaussianity. The former favors localized, event-like features;
the latter favors contextual, constraint-tracking features. Both are
measuring the activation manifold, but from different angular positions.

This suggests that the natural coordinate system recovered by ICA captures
a significant fraction of the functionally important structure in activation
space at much lower computational cost (no gradient-based dictionary
training, no hyperparameter search over sparsity levels, no per-layer
retraining). From the admissibility-field perspective, this is the expected
result: ICA is finding the eigendirections of $\nabla \fA$, which are the
natural coordinate axes of the constraint landscape; SAEs are learning a
trained dictionary that approximates those axes from data, with the
additional capacity to resolve finer-grained distinctions within each
natural cleavage plane.

\subsection{Fourth Structural Observation}

The fourth structural observation: activation space contains a latent
coordinate system prior to any trained decoder. ICA recovers this coordinate
system by finding statistically exceptional directions in the pre-existing
activation distribution. These directions correspond to constraint-regime
boundaries in the admissibility field $\fA$, and their effective receptive
fields measure the correlation length of the corresponding constraint modes.
The layer-wise shift from small to large ERF corresponds to the progressive
integration of short-range into long-range constraint fields.

%----------------------------------------------------------------------
\section{WMSD: Velocity Fields Over Possibility Space}
%----------------------------------------------------------------------

\subsection{The Teacher-Student Asymmetry}

World Model Self-Distillation \cite{stapf2026wmsd} addresses a different but
structurally related problem. The setting is instruction-conditioned video
world modeling: given an unlabeled scene image and a short task prompt, a
model must generate a plausible video trajectory showing successful task
execution. The paper evaluates on general tasks drawn from video-game
environments and real-world scenes (the WorldTasks-Bench benchmark) and
additionally transfers to robotic manipulation tasks on the DreamGen
benchmark. The examples of tasks include cutting vegetables, stacking
objects, navigating environments, and object interactions across
first-person, human-character, and robot agent settings. The training
framework involves two agents with different information access:

The \emph{Demonstrator} (teacher) receives rich conditioning:
$c_D = (\mathcal{I}, \mathcal{D})$, where $\mathcal{I}$ is the initial
observation (a scene image) and $\mathcal{D}$ is a detailed step-by-step
solution description generated by a VLM showing exactly how the task is
to be performed.

The \emph{Executor} (student) receives impoverished conditioning:
$c_E = (\mathcal{I}, \mathcal{T})$, where $\mathcal{T}$ is a short
task prompt. The Executor sees the same initial image but receives no
procedural guidance; it knows only the goal, not the path.

A concrete example drawn from the paper: for the task ``cut the carrots,''
the Demonstrator receives a detailed solution description --- ``walk to the
counter, pick up the knife, grasp the carrot, slice repeatedly, place
the pieces in the bowl'' --- generated by a VLM given the initial scene
image. The Executor receives only ``cut the carrots'' as the task prompt
together with the same initial frame. The informational asymmetry is
enormous, and yet the experimental results show the Executor eventually
surpassing the Demonstrator under VLM-based evaluation.

\subsection{What the Student Must Learn}

The training objective for the Executor is not to imitate the Demonstrator's
trajectories directly. Instead, the Executor's video world model is trained
to generate plausible future trajectories from the current state given the
goal, and a reinforcement learning signal from a vision-language model (VLM)
provides feedback on whether the generated trajectories are successful.

The critical point is what the Executor must represent in order to perform
well. It cannot memorize a trajectory $\tau^*$, because it does not have
access to the demonstration. It must instead learn the \emph{region of
trajectory space} that successfully achieves the goal --- the set of
trajectories that the Executor's generator could produce that would score
positively under the VLM reward. This is not a single path through the
workspace; it is a manifold of admissible paths.

In the language of the reachability framework, the Executor must learn an
approximation to the set $\reachset{x}$ of futures reachable from the current
state that are compatible with the task goal. This is exactly the reachability
operator from Section~4, now instantiated in the concrete setting of
instruction-conditioned video generation.

\subsection{The Velocity Field Formalism}

The WMSD architecture models trajectories using conditional flow-matching
video generators \cite{stapf2026wmsd}. The latent video state $x_t \in
\R^d$ evolves from a Gaussian base distribution $x_0 \sim p_0$ at flow
time $t = 0$ to the generated video latent $x_1 \sim p_1$ at $t = 1$.
Concretely, the Executor generates a future trajectory by integrating a
learned velocity field:
\begin{equation}
  \frac{d x_t}{dt} = v_\theta(x_t, t \mid c_E),
\end{equation}
where $x_t$ is the state at time $t \in [0,1]$, $v_\theta$ is the learned
velocity field parameterized by $\theta$, and $c_E$ is the Executor's
conditioning. Similarly, the Demonstrator's velocity field is
\begin{equation}
  \frac{d y_t}{dt} = v_{\theta'}(y_t, t \mid c_D).
\end{equation}
The fundamental object in this architecture is not a state representation but
a vector field over state space. The ``world model'' is, in its mathematical
essence, a velocity field $v(x,t)$ that governs how the state $x$ evolves
over time $t \in [0,1]$.

This is a significant ontological choice. Classical planning models represent
the world as a sequence of states connected by actions. The WMSD architecture
represents the world as a continuous flow through a state manifold. In the
former view, the primitive object is the state; in the latter, the primitive
object is the velocity field. This is closely parallel to the ontological
distinction in physics between particle mechanics (where the primitive object
is the trajectory) and field theory (where the primitive object is the field
itself). WMSD has, perhaps unwittingly, adopted the field-theoretic
formulation.

\subsection{Proposition 1: Local Agreement Controls Global Reachability}

The technical heart of the WMSD paper is the stability theorem relating local
velocity field agreement between Demonstrator and Executor to global trajectory
divergence. The paper distinguishes two variants of distillation training that
motivate the theorem. In \emph{off-policy} distillation, the student velocity
is matched to the teacher velocity on teacher-generated states: this is stable
but constrains the student only on the teacher's trajectories, so errors
compound during student rollouts. In \emph{on-policy} distillation, the
discrepancy is evaluated on student-generated states, addressing the
distribution shift at the cost of a more complex gradient. It is the on-policy
variant that motivates Proposition~1, and the empirical results show that
on-policy self-distillation continues improving past the point where
off-policy training saturates.

We reproduce the key derivation here and provide an interpretation
that is not present in the original paper.

\begin{theorem}[Velocity Field Stability; {\cite{stapf2026wmsd}}]
  \label{thm:wmsd_stability}
  Let $v_\theta$ and $v_{\theta'}$ be the Executor and Demonstrator velocity
  fields respectively, and suppose $v_{\theta'}$ is $L$-Lipschitz in its first
  argument. Let $x_0 = y_0$ be the shared initial state. Let
  $e_t = x_t - y_t$ be the trajectory error at time $t$. Suppose that
  \begin{equation}
    \int_0^1 \norm{v_\theta(x_t, t) - v_{\theta'}(x_t, t)}^2 \, dt \leq \varepsilon^2.
    \label{eq:local_matching}
  \end{equation}
  Then the terminal Wasserstein-2 distance satisfies
  \begin{equation}
    \WTwo\bigl(p_\theta(x_1 \mid c_E),\; p_{\theta'}(x_1 \mid c_D)\bigr)
    \leq e^L \varepsilon.
    \label{eq:wmsd_stability_result}
  \end{equation}
\end{theorem}

\begin{proof}[Derivation]
  By the triangle inequality for vector norms applied to the trajectory error,
  \begin{align}
    \frac{d}{dt}\norm{e_t}
    &\leq \norm{\frac{d}{dt}(x_t - y_t)}
     = \norm{v_\theta(x_t, t) - v_{\theta'}(y_t, t)}.
  \end{align}
  Adding and subtracting $v_{\theta'}(x_t, t)$:
  \begin{align}
    \frac{d}{dt}\norm{e_t}
    &\leq \norm{v_\theta(x_t, t) - v_{\theta'}(x_t, t)}
         + \norm{v_{\theta'}(x_t, t) - v_{\theta'}(y_t, t)} \\
    &\leq \delta_t + L\norm{e_t},
  \end{align}
  where $\delta_t = \norm{v_\theta(x_t, t) - v_{\theta'}(x_t, t)}$ is the
  local matching discrepancy and we have used the $L$-Lipschitz condition on
  $v_{\theta'}$. The Gr\"{o}nwall inequality applied to this differential
  inequality, with initial condition $\norm{e_0} = 0$ (since $x_0 = y_0$), yields
  \begin{equation}
    \norm{e_1} \leq e^L \int_0^1 \delta_t \, dt.
  \end{equation}
  By the Cauchy-Schwarz inequality,
  \begin{equation}
    \int_0^1 \delta_t \, dt
    \leq \left(\int_0^1 \delta_t^2 \, dt\right)^{1/2}
    \leq \varepsilon,
  \end{equation}
  where the last inequality uses assumption \eqref{eq:local_matching}. Since
  the trajectory distributions are determined by the velocity fields, the bound
  on $\norm{e_1}$ translates to the Wasserstein-2 bound
  \eqref{eq:wmsd_stability_result} via the transport argument developed in the
  following subsection.
\end{proof}

\subsection{From Trajectory Stability to Distributional Stability}

The Gr\"{o}nwall argument above establishes a sample-wise bound:
$\norm{x_1 - y_1} \leq e^L \varepsilon$ when $x_0 = y_0$. The step to a
Wasserstein-2 bound on the terminal distributions $p_\theta(x_1 \mid c_E)$
and $p_{\theta'}(x_1 \mid c_D)$ requires an explicit coupling argument. We
provide it here.

The probability density path $p_t$ for each flow is governed by the
\emph{continuity equation}:
\begin{equation}
  \partial_t p_t + \nabla \cdot (p_t v_t) = 0,
  \label{eq:continuity}
\end{equation}
which links the evolution of probability mass to the velocity field $v_t$.
The Wasserstein-2 distance between two distributions $\mu$ and $\nu$ on
$\R^d$ is defined via optimal transport:
\begin{equation}
  \WTwo(\mu, \nu)^2 = \inf_{\gamma \in \Pi(\mu,\nu)}
  \int_{\R^d \times \R^d} \norm{x - y}^2 \, d\gamma(x,y),
\end{equation}
where $\Pi(\mu, \nu)$ is the set of couplings with marginals $\mu$ and $\nu$.

We construct a specific coupling. Let $x_0 \sim p_0$ be the shared base
noise. Under the Executor's flow with velocity $v_\theta$, the point $x_0$
maps to $x_1^\theta$. Under the Demonstrator's flow with velocity $v_{\theta'}$,
the same $x_0$ maps to $x_1^{\theta'}$. The joint distribution of
$(x_1^\theta, x_1^{\theta'})$ under this shared-noise coupling is a valid
element of $\Pi(p_\theta(x_1 \mid c_E), p_{\theta'}(x_1 \mid c_D))$.

Using this coupling,
\begin{align}
  \WTwo\bigl(p_\theta(x_1 \mid c_E), p_{\theta'}(x_1 \mid c_D)\bigr)^2
  &\leq \mathbb{E}_{x_0 \sim p_0}\bigl[\norm{x_1^\theta - x_1^{\theta'}}^2\bigr] \\
  &= \mathbb{E}_{x_0 \sim p_0}\bigl[\norm{e_1}^2\bigr].
\end{align}
The Gr\"{o}nwall bound gives $\norm{e_1} \leq e^L \varepsilon$ for each
realization of $x_0$, so
\begin{equation}
  \WTwo\bigl(p_\theta(x_1 \mid c_E), p_{\theta'}(x_1 \mid c_D)\bigr)
  \leq \sqrt{\mathbb{E}[\norm{e_1}^2]}
  \leq e^L \varepsilon,
\end{equation}
where the last inequality uses the deterministic bound on $\norm{e_1}$.

The continuity equation \eqref{eq:continuity} plays a central structural role
here: it is the mechanism by which local velocity field agreement propagates
to global distributional agreement. Matching velocity fields means matching
the vector fields that transport probability mass; the Wasserstein distance
between the terminal distributions is therefore bounded by the integrated
mismatch in those transports. This connection between the continuity equation
and Wasserstein stability is precisely the bridge between the differential
and dynamical layers of the geometric triad established in Section~6.

The standard reading of this theorem is as a technical stability result: if
you train the Executor to match the Demonstrator's velocity field pointwise,
the Executor's terminal distribution will be close to the Demonstrator's.

Our reading is different and, we argue, more fundamental. Condition
\eqref{eq:local_matching} says that the Executor's velocity field agrees with
the Demonstrator's on the Executor's own trajectories. This is a condition on
\emph{local admissibility agreement}: at every point the Executor visits, the
direction of evolution it chooses must agree with the direction the Demonstrator
would choose. Result \eqref{eq:wmsd_stability_result} says that this local
agreement controls global trajectory distributions. In our language:

\begin{principleBox}
  \textbf{Local Admissibility Agreement Controls Global Reachability.}
  If the Executor matches the Demonstrator's velocity field on its own
  trajectories (local condition), then the Executor's terminal state
  distribution is close to the Demonstrator's (global consequence). Local
  velocity agreement induces global trajectory agreement.
\end{principleBox}

This is precisely the statement one would expect in a framework where the
fundamental object is the admissibility field $\fA$ rather than individual
states. The velocity field $v_\theta(x,t)$ encodes, at each state $x$, the
locally admissible direction of motion. Matching velocity fields means matching
the local admissibility structure; the global result follows by integration.

\subsection{Generation-Verification Asymmetry as Admissibility Filtering}

A recurring theme in the WMSD paper is the \emph{generation-verification
asymmetry}: it is much harder to generate a solution trajectory than to verify
whether a proposed trajectory achieves the goal. This asymmetry is pervasive in
computational problems \cite{pearl2009causality} --- SAT solving, theorem
proving, protein folding --- and the WMSD architecture exploits it by separating
the generation role (the Executor's world model, which proposes trajectories)
from the verification role (the VLM, which assesses whether proposed trajectories
are successful).

In the admissibility framework, this is an instantiation of admissibility
filtering. The Executor proposes a trajectory $\tau$; the VLM checks whether
$\tau$ belongs to the admissible set $\cF_{\mathrm{valid}}$. In the actual
WMSD implementation, the VLM (Qwen3.5-27B in the primary experiments) produces
a binary judgment on task completion, and the task reward is defined as the
log-probability difference:
\begin{equation}
  r_{\mathrm{task}}(\tau) = \log p_{\mathrm{VLM}}(\texttt{yes} \mid \tau)
  - \log p_{\mathrm{VLM}}(\texttt{no} \mid \tau),
\end{equation}
which encodes both the predicted outcome and the model's uncertainty.
The paper also introduces a consistency reward that penalizes violations
of physical plausibility and temporal coherence, providing a safeguard
against reward hacking (unrealistic object appearances, implausible motion).
The VLM is therefore not acting as a hard admissibility oracle but as a
\emph{soft} admissibility signal --- a noisy but useful measure of how deeply
inside the admissible set $\cF_{\mathrm{valid}}$ the proposed trajectory lies.
The combination of distillation regularization and soft VLM reward gives the
system a well-defined interior to explore rather than a sharp boundary to
exploit.

\subsection{Fifth Structural Observation}

The fifth structural observation: the WMSD architecture instantiates the
dynamics layer of the topology-differential structure-dynamics triad. Its
fundamental primitive is a velocity field $v(x,t)$ over possibility space,
not a state representation. The key stability theorem is a statement about
local admissibility agreement controlling global reachability. The
generation-verification asymmetry is a computational instance of admissibility
filtering. WMSD is not a world model in the classical state-representation
sense; it approximates an admissibility field.

%----------------------------------------------------------------------
\section{Distillation as Admissibility-Preserving Projection}
%----------------------------------------------------------------------

\subsection{The Compression Question}

The WMSD result that the Executor eventually surpasses the Demonstrator's
performance appears paradoxical. In ordinary distillation settings, the
student is expected to inherit the teacher's performance ceiling when trained
only to imitate the teacher directly, since the student has access to strictly
less information at inference time. How can the Executor, who receives only the
task prompt $\mathcal{T}$ while the Demonstrator receives the full solution
description $\mathcal{D}$, come to perform better than the Demonstrator?

The resolution lies in recognizing that the Demonstrator is not the objective.
The Demonstrator is a \emph{prior}: it biases the Executor's search toward the
manifold of successful trajectories, but the ultimate arbiter of success is the
VLM reward $r_{\mathrm{task}}$. Once the Executor has used the Demonstrator's
signal to initialize near the relevant region of trajectory space, the
reinforcement learning component can explore within that region and find
trajectories that are strictly better than any the Demonstrator can produce.

The deep question is: what structural information is preserved when the
Demonstrator's rich signal $(\mathcal{I}, \mathcal{D})$ is compressed to the
Executor's sparse signal $(\mathcal{I}, \mathcal{T})$? And why is enough
structure preserved for the Executor to succeed?

\subsection{The Task Sufficiency Projection Principle}

Define the compression map:
\begin{equation}
  \pi : (\mathcal{I}, \mathcal{D}) \;\mapsto\; (\mathcal{I}, \mathcal{T}).
\end{equation}
This map discards the detailed solution $\mathcal{D}$ and retains only the
task description $\mathcal{T}$. It is a projection from a high-information to
a low-information representation. Let $R(\mathcal{I}, \mathcal{D})$ and
$R(\mathcal{I}, \mathcal{T})$ denote the reachability operators associated
with the rich and sparse conditioning signals respectively --- the distributions
over future trajectories accessible under each conditioning.

\begin{principleBox}
  \textbf{Task Sufficiency Projection Principle.}
  The compression $\pi$ is \emph{admissibility-preserving} when
  \begin{equation}
    D\bigl(R(\mathcal{I}, \mathcal{D}),\; R(\mathcal{I}, \mathcal{T})\bigr)
    \leq \varepsilon
  \end{equation}
  for sufficiently small $\varepsilon$. The empirical success of WMSD
  demonstrates that $\pi$ is approximately admissibility-preserving for the
  tasks studied: despite the enormous informational asymmetry, the reachability
  structure is approximately preserved.
\end{principleBox}

This principle has a precise interpretation in terms of the admissibility field.
The projection $\pi$ is admissibility-preserving if the level sets of $\fA$
in the Demonstrator's conditioning space are approximately aligned with the
level sets of $\fA$ in the Executor's conditioning space. The detailed solution
$\mathcal{D}$ specifies a particular trajectory through $\fA$'s landscape;
the task specification $\mathcal{T}$ specifies only the target region of that
landscape. The projection preserves which region is the target, even though it
discards the specific path.

\subsection{The Role of the Reward in Reshaping the Geometry}

The WMSD reward decomposes as:
\begin{equation}
  R(\tau) = \lambda_{\mathrm{task}} \, r_{\mathrm{task}}(\tau)
            + \lambda_{\mathrm{distill}} \, r_{\mathrm{distill}}(\tau).
\end{equation}
The distillation reward $r_{\mathrm{distill}}$ measures how closely the
Executor's trajectory matches the Demonstrator's distribution. The task reward
$r_{\mathrm{task}}$ measures whether the trajectory achieves the goal. Their
combination has the geometric interpretation:
\begin{equation}
  \text{local prior (Demonstrator)} + \text{global direction (reward)}
  \;\Ra\; \text{admissibility-guided trajectory search}.
\end{equation}
The distillation term provides local geometry: it ensures that the Executor
stays near the manifold of trajectories the Demonstrator considers plausible,
preventing the RL component from wandering into degenerate regions of state
space. The reward term provides global direction: it reshapes the topology of
the admissibility field toward goal satisfaction, expanding the basin of
attraction around successful trajectories.

Together, these two terms define a constraint field that the Executor learns
to navigate. The Executor is not copying trajectories; it is learning the
geometry of admissible trajectory space as jointly defined by the Demonstrator's
prior and the reward's reshaping. This is why it can surpass the teacher: the
teacher's trajectory was one point in the admissible manifold; the Executor,
having learned the geometry of the entire manifold, can find better points.

\subsection{Detailed Instructions as Projection Artifacts}

A key conceptual conclusion of the Task Sufficiency Projection Principle is
that the detailed solution $\mathcal{D}$ is, in a precise sense, a
\emph{projection artifact} of the task representation. The admissibility
structure --- which trajectories lead to goal satisfaction --- is fully
captured by the task specification $\mathcal{T}$ together with the task reward,
without requiring the detailed procedural walkthrough. The Demonstrator's
solution provides a useful initialization and regularization signal, but the
information that is causally necessary for task competence is already present
in the task itself.

This is a non-trivial claim. It might have been the case that the detailed
procedural steps contain information about the task that is not recoverable
from the task description alone --- that certain subtleties of execution are
specified by the Demonstrator that the Executor could not independently
discover. The empirical result says otherwise: the constraint structure of the
task is largely determined by the task specification, and the Executor can
learn that structure through exploratory reinforcement learning guided by the
reward signal. Procedure is a projection artifact; constraint is invariant.

This is a direct parallel to the claim in the reachability framework
\cite{flyxion2026reachability} that much of the representational content of
a system can be derived from its future-reachability structure, without
requiring an explicit specification of all the intermediate steps.

%----------------------------------------------------------------------
\section{Frozen Processes and the Ontological Inversion}
%----------------------------------------------------------------------

\subsection{A Recurring Pattern}

We have now examined three independent empirical programs in detail. In each
case, we have identified the structural observation that the program is
making, connected it to the admissibility field $\fA$, and shown how the
connection illuminates both the empirical finding and the formal theory.
But there is a meta-level observation that has been implicit throughout
and that deserves explicit articulation: all three programs, entirely
independently, have performed the same ontological inversion.

By \emph{ontological inversion} we mean the shift from a
representation-primary account to a future-structure-primary account. In the
traditional view of language models and intelligent systems, the basic
ontological category is the \emph{object}: the model represents things ---
tokens, concepts, facts, states of the world --- and behavior is derived from
those representations. In the picture that emerges from the three programs we
have studied, the basic ontological category is the \emph{transformation}: the
model represents possibilities, constraints, and continuation dynamics, and
objects appear as derived categories that inherit their identity from the
transformational context in which they occur.

\subsection{The Empirical Inversion Table}

This inversion is not a metaphysical claim we are imposing on the data; it is
a reading of what the data actually shows. The following summary captures the
contrast:

The traditional view expects \emph{concepts} as the basic units of
representational organization: the network should have representations of
``dog,'' ``democracy,'' ``prime number,'' and so forth. What the empirical
programs find instead are \emph{constraint regimes}: the network organizes
around Q/A boundaries, patent headings, copyright notices --- not around the
concepts that the text happens to discuss.

The traditional view expects \emph{features} as the basic units of activation
structure: individual neurons or linear directions should correspond to
identifiable properties of the input. What ICALens finds instead are
\emph{continuation modes}: directions that are activated specifically when the
text enters a particular structural regime, not when a particular concept is
instantiated.

The traditional view expects \emph{states} as the basic units of world
modeling: the system should represent the current configuration of the
environment. What WMSD finds instead are \emph{trajectories}: the system is
organized around the space of future evolution, and states are points on
trajectories rather than the fundamental objects.

The traditional view expects \emph{objects} as the basic ontological category:
things with properties that persist through time. What all three programs find
are \emph{transformations}: structured modes of change that define the identity
of the context they apply to.

\subsection{Nouns as Frozen Verbs: A Mechanistic Reading}

The thesis that nouns are ``frozen verbs'' --- that object-concepts are derived
from process-concepts by a kind of temporal averaging or crystallization ---
has a long history in process philosophy \cite{kauffman1993origins,
wolfram2002nks}. In the context of the present paper, this thesis acquires
a precise mechanistic reading.

Consider again the Rank-2 cluster: 276 instances of the colon ``\texttt{:}''
after ``Q.'' The colon is, in standard linguistic analysis, a punctuation
symbol with a lexical identity. But the activation analysis reveals that the
network is not primarily representing ``colon-as-symbol''; it is representing
``colon-as-entrance-to-Q/A-regime'' --- a processual state defined by the
transformation it initiates rather than by the object it labels.

Similarly, the patent-heading cluster is not primarily representing
``heading-as-noun-phrase''; it is representing ``heading-as-entry-into-patent-
discourse-mode'' --- a structural transformation that reorganizes the admissibility
landscape. The copyright cluster is not representing ``copyright notice as a
type of legal statement''; it is representing the transformation into
copyright-language mode.

In each case, the noun-like entity (the lexical token, the surface syntactic
structure) is secondary to the verb-like entity (the transformation mode, the
constraint regime transition) that the network's geometry is actually tracking.
The activation manifold is not partitioned into things; it is partitioned into
families of admissible transformations.

\subsection{The Inversion as Empirical Consequence}

The crucial point is that this ontological inversion is not arrived at by
philosophical argument. It is not that we have decided on a priori grounds that
process-primary ontology is more elegant or more fundamental. It is that all
three empirical programs --- the activation graph, the ICA decomposition, and
the world model distillation --- have independently arrived at representations
that are organized around transformations rather than objects. The inversion
is forced by the data.

This gives us a stronger claim than process philosophy has previously been
able to make. The argument is no longer: ``conceptually, it is more natural to
think of nouns as derived from verbs.'' The argument is: ``the internal
geometry of large-scale neural language models and world models is, as an
empirical matter, organized primarily around transformational modes rather than
object categories. The traditional noun-primary ontology is not merely
philosophically inelegant; it is descriptively inaccurate at the level of the
network's actual computational structure.''

%----------------------------------------------------------------------
\section{The Admissibility Field Hypothesis}
%----------------------------------------------------------------------

\subsection{A Note on Ontological Status}

Before assembling the full formal synthesis, we pause for an important
qualification. The admissibility field $\fA : \cM \to \R_{\geq 0}$ has been
introduced and used throughout this paper as a theoretical organizing object.
This usage risks a misreading that we want to explicitly forestall.

$\fA$ is \emph{not} assumed to exist as a primitive scalar quantity that can
be directly measured from the model's weights or activations. Rather, it is
introduced as the minimal latent object whose topology, gradients, and induced
flows would jointly account for the observed clustering, ICA directions, and
velocity-field dynamics. The hypothesis succeeds insofar as it generates
testable predictions unavailable to purely representational accounts. If
tomorrow the evidence were better explained by a tensor field, a bundle, or
a sheaf, those objects would take $\fA$'s place; the structural correspondence
between the three empirical programs and the three layers of geometry would
remain.

The admissibility field is the simplest scalar projection of an underlying
geometric structure that we now make explicit.

\subsection{The Reachability Bundle and the Reachability Metric}

The most fundamental object in the framework is not the scalar admissibility
field but the \emph{reachability bundle}
\begin{equation}
  \mathcal{R} = \bigl\{R(x)\bigr\}_{x \in \cM},
\end{equation}
the collection of all reachability operators, one for each point on the
activation manifold. The admissibility field $\fA$ is one scalar functional
of this bundle: $\fA(x) = F(R(x))$ for some function $F$ (measure of admissible
mass, volume, entropy, etc.). Different choices of $F$ yield different scalar
fields, all of them shadows of the same underlying object $\mathcal{R}$.

This reordering is significant. The paper's argument has been:
\begin{equation}
  \mathcal{R} \;\lra\; \fA \;\lra\; \text{clusters, ICA directions, velocity fields}.
\end{equation}
But the deeper claim, more consistent with the prior work in this series, is
that $\mathcal{R}$ is the fundamental object and $\fA$ is merely one observable:
\begin{equation}
  \mathcal{R} \;\lra\; \text{topology of clusters}
  \quad\text{and}\quad
  \mathcal{R} \;\lra\; \text{differential structure of ICA}
  \quad\text{and}\quad
  \mathcal{R} \;\lra\; \text{dynamics of WMSD},
\end{equation}
with $\fA$ playing the role of a useful coordinate system on $\mathcal{R}$
rather than the fundamental object itself.

The reachability bundle also suggests a natural Riemannian structure on $\cM$.
Define the \emph{reachability volume} at $x$ as
\begin{equation}
  V(x) = \operatorname{Vol}\bigl(R(x)\bigr)
\end{equation}
for some measure on trajectory space. Then define a metric tensor:
\begin{equation}
  g_{ij}(x) = \frac{\partial^2}{\partial x_i \partial x_j}
  \log V(x),
  \label{eq:reachability_metric}
\end{equation}
the Hessian of the log-reachability-volume. This is the \emph{reachability
metric}: a natural Riemannian metric on activation space induced by the
geometry of reachable futures rather than by any representational distance.
Under this metric:

Activation clusters become regions of similar reachability volume --- level
sets and basins of $V(x)$ rather than of an externally defined validity measure.

ICA directions become principal curvature directions of the reachability
metric --- the eigenvectors of $g_{ij}$ at representative points, pointing
along the directions of most rapid variation in $V$.

WMSD velocity fields become geodesic flows or gradient flows through the
reachability metric, following the natural geometry of the reachability
landscape rather than an imposed objective.

This formulation tightens the mathematical backbone considerably: the
activation manifold becomes a Riemannian manifold with a canonical metric
derived entirely from the model's learned continuation structure, and all
three empirical programs become geometric measurements of that manifold.

\subsection{The Geometric Interpretability Correspondence}

We now state formally the structural correspondence that has been the central
architectural achievement of the paper. The following is not a theorem in the
strict sense --- its precise formulation would require additional assumptions
about smoothness and the form of $V(x)$ --- but it is the central conceptual
claim that all subsequent arguments elaborate.

\begin{principleBox}
\textbf{Geometric Interpretability Correspondence.}
Let $\cM$ be the normalized activation manifold and $\mathcal{R}$ the
reachability bundle. The three major families of interpretability measurement
correspond to the three foundational layers of differential geometry applied to
$\mathcal{R}$:
\begin{equation}
\begin{aligned}
  \text{Topological probes (Liu's graph)} &\;\;\longleftrightarrow\;\;
  \text{level-set structure of } \mathcal{R} \\
  \text{Differential probes (ICALens)} &\;\;\longleftrightarrow\;\;
  \text{gradient and curvature structure of } \mathcal{R} \\
  \text{Dynamical probes (WMSD)} &\;\;\longleftrightarrow\;\;
  \text{flow and transport structure of } \mathcal{R}
\end{aligned}
\end{equation}
These three families are measuring the same underlying geometric object ---
the reachability bundle --- from three different mathematical perspectives:
topological (which regions are connected), differential (which directions are
preferred), and dynamical (how states evolve). The three layers are
logically ordered: topology precedes differential structure, which precedes
dynamics.
\end{principleBox}

This correspondence is the most original conceptual contribution of the paper.
It survives even if the specific scalar field $\fA$ is replaced by a better
object. As long as the activation manifold has a reachability bundle, the
three measurement families will correspond to the three layers of geometry
applied to that bundle. The correspondence is structural, not dependent on
any particular choice of scalar summary.

\subsection{Formal Definition of the Admissibility Field}

We now give the most precise version of the scalar admissibility field, with
the understanding that this is one projection of the deeper reachability bundle
$\mathcal{R}$.

\begin{definition}[Admissibility Field, formal version]
  \label{def:admissibility_field_formal}
  Let $\cM \subset \Sph^{d-1}$ be the normalized activation manifold, and let
  $\Omega$ be the space of all future activation trajectories
  $\tau = (x_{t_0}, x_{t_1}, \ldots) \subset \cM$ beginning from various
  initial states. Let $\cF_{\mathrm{valid}} \subset \Omega$ be a measurable
  set of \emph{admissible} trajectories. Let $\mu_x$ be the probability measure
  on $\Omega$ induced by the model's continuation dynamics starting from
  state $x$. The \emph{admissibility field} is
  \begin{equation}
    \fA : \cM \lra \R_{\geq 0},
    \qquad
    \fA(x) = \mu_x\bigl(\cF_{\mathrm{valid}}\bigr)
             = \int_\Omega \mathbf{1}[\tau \in \cF_{\mathrm{valid}}] \, d\mu_x(\tau).
  \end{equation}
  Equivalently, $\fA(x)$ is the probability, under the model's continuation
  dynamics starting from $x$, that the resulting trajectory is admissible.
  The admissibility field is one scalar observable derived from the reachability
  bundle; the reachability volume $V(x) = \operatorname{Vol}(R(x))$ and the
  reachability entropy $H(R(x))$ are alternative observables with arguably
  closer correspondence to the empirical constraint-density observations.
\end{definition}

\subsection{The Four Measurements}

With the formal structure in hand, we can now state the claim that the three
empirical programs are different measurements of $\mathcal{R}$ (or equivalently,
of $\fA$ as its scalar proxy):

\begin{equation}
\begin{aligned}
  \text{Activation clusters} &\;\approx\; \text{path-connected components of superlevel sets }
    \bigl\{x : \fA(x) \geq c\bigr\} \\
  \text{ICA directions} &\;\approx\; \text{eigendirections of } g_{ij}(x)
    \text{ (or of } \nabla^2 \fA\text{)} \\
  \text{Effective Receptive Field} &\;\approx\; \text{correlation length of } \fA \\
  \text{WMSD velocity fields} &\;\approx\; \text{gradient flows through } V(x).
\end{aligned}
\end{equation}

\subsection{The Synthesis Progression}

The argument of this paper has proceeded in five steps. We now state them
explicitly as a logical sequence:

\textbf{Step 1.} Normalization is a canonical projection $\pi_\Sph$ from
$\R^d \setminus \{0\}$ to $\Sph^{d-1}$, collapsing the scale fiber $\R_{>0}$
and aligning the representation with the computationally relevant angular
geometry. This projection reveals latent geometric structure.

\textbf{Step 2.} The revealed structure clusters not around semantic topics
but around continuation constraints. Activations cluster together when they
occur at the entrance to the same constrained continuation regime, not when
they refer to the same objects or concepts.

\textbf{Step 3.} ICA finds statistically exceptional directions in the
pre-existing activation distribution. These directions correspond to
constraint-regime boundaries and approximate the curvature directions of the
reachability metric. They are natural cleavage planes of the activation manifold.

\textbf{Step 4.} WMSD's learned velocity fields represent the dynamics of
reachability-guided trajectory evolution. The fundamental stability result
(Theorem~\ref{thm:wmsd_stability}) says that local velocity agreement ---
formalized through the continuity equation and shared-noise coupling ---
controls global distributional agreement.

\textbf{Step 5.} All four observations --- clustering, ICA directions, ERF,
velocity fields --- are unified by the reachability bundle $\mathcal{R}$, of
which $\fA$ is one scalar projection. The admissibility field is the simplest
geometric entity capable of organizing this collection of observations; the
reachability bundle is the underlying geometric object of which it is an
observable.

\begin{equation}
  \boxed{
    \text{Features are emergent coordinates of admissibility topology.}
  }
\end{equation}

\subsection{The ERF Green's Function Conjecture}

The open question of formalizing ERF as a correlation length can be given a
more precise form as a conjectural field equation. Suppose the propagation of
constraint intensity along the context axis $s$ behaves like a massive scalar
field satisfying a screened Poisson equation:
\begin{equation}
  \bigl(\nabla_s^2 - m_k^2\bigr)\Phi_k(s) = J_k(s),
  \label{eq:screened_poisson}
\end{equation}
where $\Phi_k$ is the activation of ICA component $k$ as a function of context
position $s$, $J_k$ is the source term (the local constraint trigger), and
$m_k$ is an effective mass. The Green's function for
equation~\eqref{eq:screened_poisson} in one dimension decays as
\begin{equation}
  G_k(s, s_0) \sim e^{-m_k |s - s_0|},
\end{equation}
so the characteristic decay length is $\xi_k = 1/m_k$. We then conjecture:
\begin{equation}
  \ERF(k) \;\approx\; \xi_k = \frac{1}{m_k}.
\end{equation}

Under this conjecture, the empirical findings acquire a direct field-theoretic
interpretation. High-kurtosis components with small ERF correspond to large
effective mass $m_k \gg 0$: the constraint is strongly screened and decays
rapidly over context. Long-context components with large ERF correspond to
$m_k \approx 0$: the constraint behaves like a massless field and propagates
over long ranges. The layer-wise shift from small to large ERF components
corresponds, in renormalization group language, to the progressive appearance
of near-massless modes as short-range fluctuations are integrated out in
deeper layers.

This conjecture connects ERF to the correlation length of the admissibility
field in a precise way and makes quantitative predictions: fitting the screened
Poisson model to empirical ERF data would yield component-specific effective
masses $m_k$ whose distribution across layers and kurtosis values constitutes
a testable prediction.

\subsection{Open Formal Questions}

Beyond the Green's function conjecture, the admissibility field hypothesis
raises several additional formal questions.

The first is the \emph{computability} of $\fA$ and $V(x)$.
Definition~\ref{def:admissibility_field_formal} requires integration over the
space of all future trajectories, which is exponentially large. Efficient
estimation from finite samples and formal approximation guarantees are needed.

The second is the \emph{renormalization interpretation} of the component
spectrum. The power-law distribution of component sizes suggests proximity
to a percolation critical point. Can this be formalized as a
renormalization-group fixed point of the reachability bundle? Is there a
natural coarse-graining operation on $\cM$ under which $\mathcal{R}$ is
approximately self-similar?

The third is the \emph{curvature interpretation} of
Theorem~\ref{thm:wmsd_stability}. The stability bound involves the Lipschitz
constant $L$ of the teacher velocity field. Can this bound be tightened using
the sectional curvature of the reachability metric $g_{ij}$, replacing the
global constant $L$ with a local geometric invariant?

The fourth is the \emph{reachability compression theorem}. The RDR conjecture
claims that similar futures induce geometric proximity. A rigorous version would
derive a bound of the form:
\begin{equation}
  D\bigl(R(x), R(y)\bigr) \leq \varepsilon
  \quad\Ra\quad
  \norm{\hh_x - \hh_y} \leq f(\varepsilon)
\end{equation}
from predictive sufficiency arguments about the training objective, without
assuming it empirically. Such a theorem would make the RDR conjecture a
consequence of the statistical structure of next-token prediction training,
and would provide the foundational theorem from which all other results in
the paper could be derived.

%----------------------------------------------------------------------
\section{Predictions and Falsification Criteria}
%----------------------------------------------------------------------

A theoretical framework earns its place not only by organizing known
observations but by generating predictions that distinguish it from
alternatives. The admissibility field hypothesis and the RDR conjecture
make several empirical predictions that could in principle be tested with
existing interpretability infrastructure.

\textbf{Prediction 1: Reachability predicts activation neighborhoods better
than semantic similarity.} If the RDR conjecture is correct, then the
reachability distance $D(R(a_1), R(a_2))$ should predict activation-space
proximity $\norm{\hh_{a_1} - \hh_{a_2}}$ better than any semantic similarity
measure (topic similarity, distributional word vectors, etc.). A falsification:
systematic evidence that semantic similarity outperforms reachability distance
as a predictor of activation-space proximity would challenge the RDR conjecture
directly.

\textbf{Prediction 2: ICA directions with highest kurtosis correspond to
strongest future-space partition boundaries.} If ICA directions approximate
eigendirections of the reachability metric $g_{ij}$, then the most non-Gaussian
components should correspond to the steepest boundaries between distinct
continuation regimes. A falsification: high-kurtosis ICA directions found to
correspond to arbitrary lexical distinctions rather than continuation-regime
transitions would break the connection between non-Gaussianity and admissibility
gradient structure.

\textbf{Prediction 3: Distillation success correlates with reachability
structure preservation.} The Task Sufficiency Projection Principle predicts
that distillation succeeds when $D(R(\mathcal{I}, \mathcal{D}), R(\mathcal{I},
\mathcal{T})) \leq \varepsilon$. A testable consequence: distillation should
fail or degrade systematically when the task specification $\mathcal{T}$ loses
information about the shape of the future trajectory manifold (e.g.\ when the
task prompt is severely ambiguous) even if the representation of the goal state
is preserved. A falsification: successful distillation in cases where the
reachability structure is demonstrably not preserved.

\textbf{Prediction 4: The ERF--kurtosis outliers are the most interpretively
important components.} The observed correlation between ERF and excess kurtosis
($\rho \approx -0.5$) is moderate, not decisive. This means there exist
components with large ERF and large kurtosis simultaneously. These are
components that depend on long-range context but activate only rarely and
strongly. The admissibility framework predicts that these outliers correspond
to the most important structural boundaries in the model: discourse-mode
transitions, document-level structural boundaries, latent task switches ---
precisely the places where the admissibility landscape changes most sharply
over long contextual ranges. A systematic annotation of these outlier
components constitutes a targeted test of this prediction.

\textbf{Prediction 5: Activation clusters predict continuation similarity
better than semantic similarity.} Liu's clusters are organized around
constrained continuation regimes rather than semantic topics. A direct test:
for any two tokens in the same cluster, their empirically observed continuation
distributions should be more similar than their semantic similarity would
predict. Conversely, two semantically similar tokens from different clusters
should have more divergent continuation distributions than semantic similarity
would suggest. This prediction can be tested directly using the activation
graph from \cite{liu2026normgeometry} paired with continuation-distribution
measurements.

Once the framework generates predictions of this kind, the admissibility field
stops being merely an interpretation and becomes a research program. Each
prediction constitutes a constraint on the reachability bundle $\mathcal{R}$
and each falsification would require a structural revision of the framework.

%----------------------------------------------------------------------
\section{Conclusion}
%----------------------------------------------------------------------

\subsection{Summary of the Argument}

We began with three empirical results that appeared, on their surfaces, to be
about entirely different things: the geometry of activation vectors on a
hypersphere, the statistical structure of independent component directions, and
the distillation of step-by-step solutions into high-level task performance.
We have argued that these three results are different measurements of a single
underlying geometric object --- the reachability bundle $\mathcal{R}$ over the
activation manifold, summarized for present purposes by the scalar admissibility
field $\fA : \cM \to \R_{\geq 0}$ --- and that recognizing this common origin
explains the structural similarities among them and opens several new directions
for both theory and practice.

The central conceptual claim is the Reachability Determines Representation
conjecture: the organizational structure of the activation manifold is shaped by
future-reachability proximity classes rather than by present-state feature
categories. We have been careful throughout to distinguish what the evidence
directly shows (geometrically clustered states share similar continuation
structures) from what the conjecture asserts in the causal direction (similar
futures induce geometric proximity). Both directions are plausible given the
evidence; only the latter requires empirical verification beyond what the papers
analyzed here provide.

The central formal contribution is the Geometric Interpretability
Correspondence: the observation that the three empirical programs instantiate
the three foundational layers of differential geometry --- topology
(Liu's graph), differential structure (ICALens directions), dynamics (WMSD
velocity fields) --- applied to the same underlying geometric object. This
correspondence is structurally robust: it holds even if the specific scalar
field $\fA$ is replaced by a more refined object such as the reachability
volume $V(x)$ or the full reachability bundle $\mathcal{R}$.

A secondary formal contribution is the clarification that the $\varepsilon$-future
proximity relation $\approx_R^\varepsilon$ is a uniform tolerance structure
rather than an equivalence relation for $\varepsilon > 0$, and the addition of
an explicit continuity-equation derivation bridging the sample-wise Gr\"{o}nwall
bound to the Wasserstein-2 distributional stability result of
Theorem~\ref{thm:wmsd_stability}.

\subsection{Implications for Interpretability}

If the admissibility field hypothesis is correct, it has significant
implications for the agenda of mechanistic interpretability. The standard
agenda assumes that the natural decomposition of activation space is
\emph{featural}: the goal is to identify a set of directions or sparse
features such that each activation can be expressed as a combination of
interpretable feature-values. The dominant tools --- sparse autoencoders,
circuit analysis, feature visualization --- are all designed to find and
analyze these featural components.

Our argument suggests that this featural decomposition, while useful, may be
an \emph{emergent coordinate system} over a more fundamental topological
structure. Features are not the primitive objects; they are the natural
coordinates of the admissibility landscape. The primary structure is the
topology of the level sets of $\fA$ --- which regions are connected, which
are isolated, which are on the boundary between regimes --- and features
appear as the differential structure overlaid on this topology.

This suggests a different methodological priority: understand the topology of
$\fA$ first, then derive the feature coordinates from the topology. Liu's
activation graph is already doing this: it is recovering the topology directly,
without going through features. ICALens is recovering the differential
structure, one level more refined than the topology but still more fundamental
than a trained feature dictionary. A fully topological interpretability program
would begin with the level-set structure of $\fA$ and derive features as its
natural coordinate system.

\subsection{Implications for Alignment}

The alignment implications are correspondingly reframed. Standard approaches
to alignment focus on specifying rewards, preferences, and values, then
training systems to optimize for them. This is the representational paradigm:
we specify what we want (a representation of desired behavior) and train the
system to match it.

If the internal geometry of intelligent systems is organized around
admissibility fields rather than reward functions, the natural alignment
target is not a reward specification but a \emph{geometry of admissible
futures}. We do not want to tell the system what is good; we want to shape the
constraint field within which the system moves so that the admissible future
set coincides with the futures we prefer. This is a geometric problem, not an
optimization problem.

More precisely: alignment may depend less on specifying rewards and more on
shaping the geometry of admissible futures. The technical challenge is to
understand $\fA$ well enough to know what changes to training, architecture,
or data would deform its level sets in the desired directions.

\subsection{The Broader Convergence}

A final word on the significance of the convergence itself. In the history of
science, convergences of this kind --- where independent empirical programs
arrive at the same formal structure from different directions --- are often
signals that the formal structure in question is genuinely real. The periodic
table was confirmed not by a single experiment but by the convergence of
electrochemistry, spectroscopy, and atomic mass measurements. General
relativity was confirmed by the convergence of perihelion precession, light
deflection, and gravitational redshift. When independent observers, using
independent methods, find the same geometric structure, the most parsimonious
conclusion is that the structure is there.

We are not claiming that the admissibility field has been confirmed with the
certainty of general relativity. But the convergence of three independent
research programs on the same geometric picture --- activation space organized
around continuation constraints rather than object categories, structured
by level sets, gradient directions, correlation lengths, and flow fields of
a single admissibility function --- is at minimum a strong suggestion that
the right language for activation space is geometric before it is featural,
topological before it is categorical, and dynamical before it is static.

The features are not wrong. They are just not fundamental. They are emergent
coordinates of a constraint topology that the network has been learning,
through billions of prediction tasks, to navigate. The activation manifold is
a map of admissible futures. Interpretability, on this view, is not the
project of reading off what the map says about the present moment. It is the
project of understanding the geometry of the territory the map describes.

%----------------------------------------------------------------------
\begin{thebibliography}{99}

\bibitem{amari2016information}
S.~Amari.
\newblock {\em Information Geometry and Its Applications}.
\newblock Springer, 2016.

\bibitem{absil2008optimization}
P.-A.~Absil, R.~Mahony, and R.~Sepulchre.
\newblock {\em Optimization Algorithms on Matrix Manifolds}.
\newblock Princeton University Press, 2008.

\bibitem{barabasi2016network}
A.-L.~Barab\'{a}si.
\newblock {\em Network Science}.
\newblock Cambridge University Press, 2016.

\bibitem{bricken2023monosemantic}
T.~Bricken et al.
\newblock Towards monosemanticity: Decomposing language models with dictionary
  learning.
\newblock {\em Transformer Circuits Thread}, Anthropic, 2023.

\bibitem{coifman2006diffusion}
R.~Coifman and S.~Lafon.
\newblock Diffusion maps.
\newblock {\em Applied and Computational Harmonic Analysis},
  21(1):5--30, 2006.

\bibitem{nadler2006diffusion}
B.~Nadler, S.~Lafon, R.~Coifman, and I.~Kevrekidis.
\newblock Diffusion maps, spectral clustering and reaction coordinates of
  dynamical systems.
\newblock {\em Applied and Computational Harmonic Analysis}, 21(1):113--127,
  2006.

\bibitem{cunningham2023sae}
L.~Cunningham et al.
\newblock Sparse autoencoders find highly interpretable features in language
  models.
\newblock {\em arXiv preprint arXiv:2309.08600}, 2023.

\bibitem{elhage2022toy}
N.~Elhage et al.
\newblock Toy models of superposition.
\newblock {\em Transformer Circuits Thread}, Anthropic, 2022.

\bibitem{flyxion2026admissibility}
Flyxion.
\newblock {\em The Admissibility Field}.
\newblock Zenodo, 2026.

\bibitem{flyxion2026frozen}
Flyxion.
\newblock {\em Frozen Processes: Reachability Ontology and the Geometry of
  Transformation}.
\newblock Zenodo, 2026.

\bibitem{flyxion2026hidden}
Flyxion.
\newblock {\em Hidden Manifolds: Projection, Sufficiency, and Reachability in
  High-Dimensional Systems}.
\newblock Zenodo, 2026.

\bibitem{flyxion2026reachability}
Flyxion.
\newblock {\em From Preservation to Reachability: Semantic Continuity in an
  Age of Drift}.
\newblock Zenodo, 2026.

\bibitem{friston2010free}
K.~Friston.
\newblock The free-energy principle: A unified brain theory?
\newblock {\em Nature Reviews Neuroscience}, 11(2):127--138, 2010.

\bibitem{ha2018worldmodels}
D.~Ha and J.~Schmidhuber.
\newblock World models.
\newblock {\em arXiv preprint arXiv:1803.10122}, 2018.

\bibitem{hassabis2024agents}
D.~Hassabis and D.~Silver.
\newblock From world models to generalist agents.
\newblock {\em Communications of the ACM}, 2024.

\bibitem{holland1992adaptation}
J.~Holland.
\newblock {\em Adaptation in Natural and Artificial Systems}.
\newblock MIT Press, 1992.

\bibitem{hyvarinen2000ica}
A.~Hyv\"{a}rinen and E.~Oja.
\newblock Independent component analysis: Algorithms and applications.
\newblock {\em Neural Networks}, 13(4--5):411--430, 2000.

\bibitem{hyvarinen2001book}
A.~Hyv\"{a}rinen, J.~Karhunen, and E.~Oja.
\newblock {\em Independent Component Analysis}.
\newblock Wiley, 2001.

\bibitem{kauffman1993origins}
S.~Kauffman.
\newblock {\em The Origins of Order}.
\newblock Oxford University Press, 1993.

\bibitem{lecun2022path}
Y.~LeCun.
\newblock A path towards autonomous machine intelligence.
\newblock Open Review, 2022.

\bibitem{lee2013smooth}
J.~M.~Lee.
\newblock {\em Introduction to Smooth Manifolds}.
\newblock Springer, 2013.

\bibitem{liu2026normgeometry}
S.~Liu.
\newblock Normalization makes activation geometry imaginable.
\newblock Research Notes, June~9, 2026.
\newblock \url{https://liusida.com/research-notes/norm_act_geometry/}

\bibitem{liu2026icalens}
S.~Liu and F.~Han.
\newblock {ICA} lens: Interpreting language models without training another
  dictionary.
\newblock {\em arXiv preprint arXiv:2606.11722}, 2026.

\bibitem{newman2018networks}
M.~Newman.
\newblock {\em Networks}.
\newblock Oxford University Press, 2018.

\bibitem{olah2020circuits}
C.~Olah et al.
\newblock Zoom in: An introduction to circuits.
\newblock {\em Distill}, 2020.

\bibitem{pearl2009causality}
J.~Pearl.
\newblock {\em Causality}.
\newblock Cambridge University Press, 2009.

\bibitem{roweis2000lle}
S.~Roweis and L.~Saul.
\newblock Nonlinear dimensionality reduction by locally linear embedding.
\newblock {\em Science}, 290(5500):2323--2326, 2000.

\bibitem{rumelhart1986representations}
D.~Rumelhart, G.~Hinton, and R.~Williams.
\newblock Learning internal representations by error propagation.
\newblock In {\em Parallel Distributed Processing}, 1986.

\bibitem{silver2021reward}
D.~Silver et al.
\newblock Reward is enough.
\newblock {\em Artificial Intelligence}, 299, 2021.

\bibitem{stapf2026wmsd}
S.~Stapf, P.~Acuaviva~Huertos, A.~Davtyan, and P.~Favaro.
\newblock World model self-distillation: Training world models to solve general
  tasks.
\newblock {\em arXiv preprint arXiv:2606.12072}, 2026.
\newblock University of Bern.

\bibitem{stauffer1994percolation}
D.~Stauffer and A.~Aharony.
\newblock {\em Introduction to Percolation Theory}.
\newblock Taylor \& Francis, 1994.

\bibitem{tenenbaum2000isomap}
J.~Tenenbaum, V.~de~Silva, and J.~Langford.
\newblock A global geometric framework for nonlinear dimensionality reduction.
\newblock {\em Science}, 290(5500):2319--2323, 2000.

\bibitem{wolfram2002nks}
S.~Wolfram.
\newblock {\em A New Kind of Science}.
\newblock Wolfram Media, 2002.

\bibitem{zhang2019rmsnorm}
B.~Zhang and R.~Sennrich.
\newblock Root mean square layer normalization.
\newblock {\em arXiv preprint arXiv:1910.07467}, 2019.

\end{thebibliography}

\end{document}
